\chapter{Imagen determinantal y composición hacia Birch--Swinnerton--Dyer}

\section{Unidad genealógica y producto fibrado}

Sea $E/\mathbb Q$ una curva elíptica y sea
$\mathfrak G_{\rm det}(E)$ la imagen determinantal autónoma completa. Se
denota por $\mathcal T_m^{\rm surv}$ su categoría de historias
supervivientes. El propietario determinantal precedente construye desde una
unidad común la cadena de líneas
\eqref{eq:amp-bsd-linea-determinante}--\eqref{eq:amp-bsd-cadena-lineas}, el
comparador Kato--FERS
\eqref{eq:amp-bsd-phi-kato-fers}--\eqref{eq:amp-bsd-inversa-kato-fers}, el
canal Poitou--Tate
\eqref{eq:amp-bsd-comparador-pt-q}--\eqref{eq:amp-bsd-comparador-pt} y el
pegado local--global
\eqref{eq:amp-bsd-triangulo-local}--\eqref{eq:amp-bsd-mapa-localizacion}.
La tupla de \eqref{eq:pdf2-g-det-explicito} y las identidades de
\eqref{eq:pdf2-g-det-identidades-dato} son la publicación conjunta de esa
construcción upstream: no se aceptan como axiomas terminales ni se eligen a
partir del rango o del término principal. Este capítulo conserva y compone
esas flechas hasta la conclusión de Birch--Swinnerton--Dyer.

El objeto de coeficientes no es sólo la fibra $3$-primaria. Se fija primero
el complejo libre integral basado
\begin{equation}\label{eq:pdf2-bsd-gmz}
 G_{m,\mathbb Z}:=N_*\bigl(N\mathcal T_m^{\rm surv};\mathbb Z\bigr),
\end{equation}
y, para cada primo $\ell$ y cada $n\geq1$, se ponen
\begin{equation}\label{eq:pdf2-bsd-coeficientes-todos-ell}
 A_{\ell,n}:=E[\ell^n](\overline{\mathbb Q}),\qquad
 T_\ell E:=\varprojlim_nA_{\ell,n},\qquad
 V_\ell E:=T_\ell E\otimes_{\mathbb Z_\ell}\mathbb Q_\ell,
 \qquad
 G_{m,\ell,n}:=G_{m,\mathbb Z}\otimes_{\mathbb Z}^{\mathbf L}A_{\ell,n}.
\end{equation}
Las transiciones no quedan implícitas: son
\begin{equation}\label{eq:pdf2-bsd-transiciones-ell-n}
 q_{\ell,n}:A_{\ell,n+1}\twoheadrightarrow A_{\ell,n},
 \quad x\longmapsto\ell x,
 \qquad
 Q_{\ell,n}:=1\otimes q_{\ell,n}:G_{m,\ell,n+1}\to G_{m,\ell,n},
\end{equation}
La transición de coeficientes en el complejo de Manin, junto con
$Q_{\ell,n}$, induce
$\widetilde Q_{\ell,n}:P_{E,\ell,n+1}^{\rm gen}\to
P_{E,\ell,n}^{\rm gen}$. Estos mapas satisfacen
$Q_{\ell,n}d=dQ_{\ell,n}$,
$(Q_{\ell,n}\otimes Q_{\ell,n})\Delta_{\rm AW}
=\Delta_{\rm AW}Q_{\ell,n}$ y
$\widetilde Q_{\ell,n}\widehat\Pi_{0,\ell,n+1}
=\widehat\Pi_{0,\ell,n}Q_{\ell,n}$. Así el límite que define $T_\ell E$
se toma dentro de la misma familia de complejos, no sólo en los grupos de
coeficientes.
La antigua notación $A_n$ designa únicamente
$A_{3,n}$; no gobierna la publicación aritmética. Las matrices de frontera,
Alexander--Whitney, orientación y carry están escritas sobre $\mathbb Z$ en
la base de historias. Por ello, para
$R\in\{\mathbb Z_\ell,\mathbb Z/\ell^n,\mathbb Q_\ell\}$,
\begin{equation}\label{eq:pdf2-bsd-base-change-integral}
 G_{m,R}=G_{m,\mathbb Z}\otimes_{\mathbb Z}^{\mathbf L}R,
 \qquad
 d_{m,R}=d_{m,\mathbb Z}\otimes1,
 \qquad
 \Delta_{{\rm AW},R}=\Delta_{{\rm AW},\mathbb Z}\otimes1.
\end{equation}
Así, el paso de $3$ a otro primo no es una extensión inexistente
$\mathbb Z_3\to\mathbb Z_\ell$: todas las fibras se obtienen por cambio de
base del mismo propietario integral.

Para abreviar, las fórmulas siguientes se escriben sobre un coeficiente
$A_{\ell,n}$ fijo y son compatibles con $n$, con $\ell$ y con la base
integral. El complejo normalizado
\begin{equation}\label{eq:pdf2-bsd-gmn}
 G_{m,\ell,n}:=N_*\bigl(N\mathcal T_m^{\rm surv};A_{\ell,n}\bigr)
\end{equation}
conserva la diagonal Alexander--Whitney
\begin{equation}\label{eq:pdf2-bsd-aw}
 \Delta_{\rm AW}[x_0\to\cdots\to x_q]
 =\sum_{j=0}^q[x_0\to\cdots\to x_j]\otimes
                [x_j\to\cdots\to x_q]
\end{equation}
y la counidad
$\varepsilon_{\rm AW}[x]=1$ en vértices, cero en grados positivos. La historia
se acopla al complejo de Manin mediante
\begin{equation}\label{eq:pdf2-bsd-pullback}
 P_{E,\ell,n}^{\rm gen}:=
 P_{E,\ell,n}^{\rm Man}\times_{A_{\ell,n}[0]}G_{m,\ell,n}.
\end{equation}
Si $s_{\ell,n}:A_{\ell,n}[0]\to P_{E,\ell,n}^{\rm Man}$ es la sección
basada, entonces
\[
 \widehat\Pi_{0,\ell,n}(c)=
 (s_{\ell,n}\varepsilon_{\rm AW}(c),c)
\]
conserva simultáneamente la coordenada modular y la historia. Esta parte de la
construcción es material y no usa el término principal de $L(E,s)$. Los
cuadrados
\begin{equation}\label{eq:pdf2-bsd-pullback-base-change}
 P_{E,\mathbb Z}^{\rm gen}\otimes_{\mathbb Z}^{\mathbf L}
      \mathbb Z/\ell^n
 \xrightarrow{\ \sim\ }P_{E,\ell,n}^{\rm gen},
 \qquad
 \widehat\Pi_{0,\mathbb Z}\otimes1=\widehat\Pi_{0,\ell,n}
\end{equation}
prueban que unidad, pullback y counidad son los mismos para todos los
coeficientes.

Las dos publicaciones integrales y sus cambios de base
\begin{equation}\label{eq:pdf2-bsd-copub}
 P_{E,R}=(P_{E,R}^K,P_{E,R}^{\rm PT}),
 \qquad P_{E,R}^K(1_E)=z_{K,R},\quad
 P_{E,R}^{\rm PT}(1_E)=z_{{\rm PT},R}
\end{equation}
deben proceder de la misma unidad $1_E$. La dualidad reducida se tipa por
\[
 D_3(M):=\operatorname{RHom}_\Lambda(M^\iota,\Lambda)[-3],
 \qquad
 X^{\rm orth}:C_E^{\rm red}\xrightarrow{\sim}D_3(C_E^{\rm red}).
\]
Con $L_{{\rm root},\ell,n}=A_{\ell,n}z_{{\rm PT},\ell,n}$ y el proyector de complejos
$p_{\rm root}$, la normalización basada da
\begin{equation}\label{eq:pdf2-bsd-raiz}
 p_{\rm root}z_{K,\ell,n}=z_{{\rm PT},\ell,n}.
\end{equation}

\section{Filler local y descomposición global sin borrar el retículo finito}

La descomposición se escribe primero sobre el propietario integral. Para
cada etiqueta finita--singular $\lambda=(\ell,m,n,\ldots)$ existe un
coeficiente de incidencia $a_\lambda=\ell c_\lambda\in\mathbb Z$ fijado por
$\operatorname{Rel}_{\rm det}$; en la fibra correspondiente
$A_{\ell,n}$ actúa por reducción de ese mismo entero. El bloque libre
integral con $g$ primitivo es
\begin{equation}\label{eq:pdf2-bsd-bloque-normal}
 \partial f=a_\lambda e+b,\qquad
 \partial e=g,\qquad
 \partial b=-a_\lambda g.
\end{equation}
Tensorizarlo con la fibra indicada por $\lambda$ devuelve literalmente
$a_\lambda=\ell c_\lambda$. En consecuencia, si
$z_K=ar+ue+vb$ es ciclo y $z_{\rm PT}=r$, entonces
$\partial z_K=(u-a_\lambda v)g=0$ implica $u=a_\lambda v$. La igualdad raíz
\eqref{eq:pdf2-bsd-raiz} da $a=1$ y, por tanto,
\begin{equation}\label{eq:pdf2-bsd-filler-local}
 h_*:=-vf,\qquad \partial h_*=z_{\rm PT}-z_K.
\end{equation}

La multisección $\Sigma_{\rm det}^{\rm mult}$ etiqueta cada generador por
\begin{equation}\label{eq:pdf2-bsd-etiqueta-completa}
 \lambda=(\ell,m,n,\gamma,\operatorname{or},\operatorname{Carr},
 \operatorname{Mem},\partial,\tau^{\rm term},\operatorname{loc}).
\end{equation}
El orden basado separa tres clases disjuntas y exhaustivas: la unidad raíz,
los bloques finita--singular con terminal $h_*$ y las celdas del retículo
integral local--global. Como $\partial_{\rm det}$ conserva todas las
etiquetas salvo el descenso prescrito por carry, estas clases definen
proyectores de complejos $p_{\rm root}$, $p_{\rm FS}$ y $p_{\rm lat}$ con
\begin{equation}\label{eq:pdf2-bsd-tres-proyectores}
 p_{\rm root}+p_{\rm FS}+p_{\rm lat}=I,
 \qquad p_ip_j=\delta_{ij}p_i.
\end{equation}

\begin{theorem}[Descomposición global basada]
\label{thm:pdf2-bsd-descomposicion-global}
Existe un isomorfismo de complejos compatible con lugares, orientación,
corestricción y dualidad
\begin{equation}\label{eq:pdf2-bsd-descomposicion-global}
 \Theta_E:\mathcal C_E^{\rm loc,red}
 \xrightarrow{\ \sim\ }
 \mathcal C_E^{\rm FS,ctr}\oplus\mathcal L_E,
\end{equation}
donde
\begin{equation}\label{eq:pdf2-bsd-sumando-fs-y-lattice}
 \begin{aligned}
 \mathcal C_E^{\rm FS,ctr}
  &:=\bigoplus_{\lambda\in\Lambda_E^{\rm FS}}
   [\mathbb Zf_\lambda\to
    \mathbb Ze_\lambda\oplus\mathbb Zb_\lambda
                    \to\mathbb Zg_\lambda],\\
 \mathcal L_E&:=[M_E\xrightarrow{D_E}N_E].
 \end{aligned}
\end{equation}
Cada bloque del primer sumando tiene el diferencial integral
\eqref{eq:pdf2-bsd-bloque-normal}; su fibra $(\ell,n)$ se obtiene por
$-\otimes_{\mathbb Z}^{\mathbf L}A_{\ell,n}$. El complejo $\mathcal L_E$
conserva, sin
contraerlos, los divisores de Smith y las componentes
Tate--Shafarevich--Tamagawa.
\end{theorem}

\begin{proof}
Se ordenan los generadores por la etiqueta
\eqref{eq:pdf2-bsd-etiqueta-completa}. La unidad común ocupa el primer
bloque y se elimina por $q_{\rm root}$. Para cada terminal
$(z_{\rm PT},h_*)$ las relaciones de $\operatorname{Rel}_{\rm det}$
identifican exactamente cuatro generadores
$(f_\lambda,e_\lambda,b_\lambda,g_\lambda)$; la matriz de incidencia es
\eqref{eq:pdf2-bsd-bloque-normal}, con $g_\lambda$ primitivo porque su
etiqueta de llegada aparece una sola vez en la base orientada. Las celdas
restantes tienen etiqueta local--global y forman las dos redes libres
$M_E,N_E$. No hay cuarta clase: la multisección es completa y
$\operatorname{Surv}_{\rm det}$ exige a cada historia o bien un terminal
FS o bien una coordenada de red. El cambio de bases es unitriangular con
diagonal uno, pues una frontera sólo introduce etiquetas anteriores. Esto
construye $\Theta_E$ y prueba simultáneamente su exhaustividad y su
compatibilidad con las bases. La dualidad $X^{\rm orth}$ intercambia las
dos redes y conserva cada bloque FS, de modo que también entrelaza
$\Theta_E$.
\end{proof}

Sobre un bloque normal defínase
\begin{equation}\label{eq:pdf2-bsd-hfs-generadores}
 H_\lambda(g_\lambda)=e_\lambda,
 \qquad H_\lambda(b_\lambda)=f_\lambda,
 \qquad H_\lambda(e_\lambda)=H_\lambda(f_\lambda)=0.
\end{equation}
Un cálculo sobre los cuatro generadores da
$\partial H_\lambda+H_\lambda\partial=I$. Por tanto
\begin{equation}\label{eq:pdf2-bsd-hfs-global}
 H_{\rm FS}:=\Theta_E^{-1}
 \left(\bigoplus_{\lambda\in\Lambda_E^{\rm FS}}H_\lambda\oplus0\right)
 \Theta_E,
 \qquad
 \partial H_{\rm FS}+H_{\rm FS}\partial=p_{\rm FS}.
\end{equation}
La coordenada terminal de $\mathfrak G_{\rm det}(E)$ asegura
$q_{\rm root}z_K\in\operatorname{im}p_{\rm FS}$; entonces
$h_*=-H_{\rm FS}(q_{\rm root}z_K)$ reproduce
\eqref{eq:pdf2-bsd-filler-local}. Es esencial que el lado derecho de
\eqref{eq:pdf2-bsd-hfs-global} sea $p_{\rm FS}$, no $q_{\rm root}$:
contraer el retículo $\mathcal L_E$ borraría los factores finitos que BSD
debe publicar.

\section{Propietario integral y cambio de base a todos los primos}

La descomposición anterior se ejecuta primero en la base integral de
historias, no en una fibra primaria. Para
\[
 R\in\{\mathbb Z,\mathbb Z_\ell,\mathbb Z/\ell^n,
          \mathbb Q_\ell\}
\]
se fijan
\begin{equation}\label{eq:pdf2-bsd-complejos-base-change}
 \mathcal C_{E,R}^{\rm loc,red}:=
 \mathcal C_{E,\mathbb Z}^{\rm loc,red}
       \otimes_{\mathbb Z}^{\mathbf L}R,
 \qquad
 \mathcal L_{E,R}:=
 [M_E\otimes R\xrightarrow{D_E\otimes1}N_E\otimes R].
\end{equation}
El orden unitriangular de las etiquetas usado en
\cref{thm:pdf2-bsd-descomposicion-global} tiene coeficientes enteros. En
consecuencia,
\begin{equation}\label{eq:pdf2-bsd-theta-base-change}
 \Theta_{E,R}=\Theta_{E,\mathbb Z}\otimes1,
 \qquad
 H_{{\rm FS},R}=H_{{\rm FS},\mathbb Z}\otimes1,
 \qquad
 \partial H_{{\rm FS},R}+H_{{\rm FS},R}\partial=p_{{\rm FS},R}.
\end{equation}
En particular, la fibra $3$-primaria es una de estas especializaciones y no
puede ocultar las fibras $\ell\ne3$.

Antes de formar un cokernel finito se construye el inverso racional. El orden
de terminales de $\Sigma_{\rm det}^{\rm mult}$ da bases
$(m_1,\ldots,m_t)$ de $M_E$ y $(n_1,\ldots,n_t)$ de $N_E$ y una biyección
basada $\sigma_E$ entre ambas listas. Si
\[
 d_j:=\bigl\langle n_{\sigma_E(j)}^\vee,D_Em_j\bigr\rangle,
\]
la relación terminal y la dualidad $X^{\rm orth}$ dan $d_j\ne0$ antes de
tomar dimensiones. Defínanse
\begin{equation}\label{eq:pdf2-bsd-d0-nilpotente}
 \begin{aligned}
  D_{E,0}:M_E\otimes\mathbb Q&\longrightarrow N_E\otimes\mathbb Q,
  &D_{E,0}m_j&=d_jn_{\sigma_E(j)},\\
  N_E^{\rm tri}&:=D_{E,0}^{-1}(D_E\otimes\mathbb Q)-I.
 \end{aligned}
\end{equation}
Una frontera sólo contiene su terminal y etiquetas estrictamente anteriores;
por ello $N_E^{\rm tri}$ es estrictamente triangular en la lista ordenada y
$(N_E^{\rm tri})^t=0$. El operador
\begin{equation}\label{eq:pdf2-bsd-j-e-inverso-racional}
 \boxed{
 J_E:=\left(\sum_{r=0}^{t-1}(-N_E^{\rm tri})^r\right)D_{E,0}^{-1}:
 N_E\otimes\mathbb Q\longrightarrow M_E\otimes\mathbb Q }
\end{equation}
se calcula, por tanto, con un número finito de operaciones sobre la matriz de
incidencia. Como
$D_E\otimes\mathbb Q=D_{E,0}(I+N_E^{\rm tri})$, la suma geométrica finita
prueba las dos identidades
\begin{equation}\label{eq:pdf2-bsd-j-e-dos-lados}
 J_E(D_E\otimes\mathbb Q)=I_{M_E\otimes\mathbb Q},
 \qquad
 (D_E\otimes\mathbb Q)J_E=I_{N_E\otimes\mathbb Q}.
\end{equation}
Ésta es la contracción racional del retículo residual; no usa rango,
$\operatorname{Sha}(E)$, órdenes de grupos ni el término principal.

La matriz integral $D_E$ conserva todos sus divisores elementales. Para cada
primo se define, sin inferirlo desde la fibra $3$-primaria,
\begin{equation}\label{eq:pdf2-bsd-f-ell-def}
 D_{E,\ell}:=D_E\otimes\mathbb Z_\ell,
 \qquad
 F_{E,\ell}:=\operatorname{coker}D_{E,\ell}.
\end{equation}
Las dos identidades de \eqref{eq:pdf2-bsd-j-e-dos-lados} implican que
$F_E:=\operatorname{coker}D_E$ es finito. Por tanto su
descomposición primaria es finita y exacta:
\begin{equation}\label{eq:pdf2-bsd-f-descomposicion-primaria}
 F_E=\operatorname{coker}D_E
 \xrightarrow{\ \sim\ }
 \bigoplus_{\ell}F_{E,\ell},
 \qquad
 [n]\longmapsto([n\otimes1]_\ell)_\ell.
\end{equation}
La inversa se obtiene generador a generador por el teorema chino del resto
aplicado a cada divisor de Smith. Esta es la arista que permite publicar
simultáneamente $\operatorname{Sha}(E)$ y todos los factores de Tamagawa.

\section{Bocksteins de todas las páginas}

Sea
\begin{equation}\label{eq:pdf2-bsd-s-universal}
 S_{E,0}(T):V_0[[T]]\longrightarrow W_0[[T]]
\end{equation}
el diferencial basado racional construido por las matrices integrales del
lector determinantal. Para cada primo $\ell$ y para la carta compleja se
obtiene por cambio de base
\begin{equation}\label{eq:pdf2-bsd-s-todos-ell}
 S_{E,\ell}(T):=S_{E,0}(T)\otimes_{\mathbb Q}\mathbb Q_\ell,
 \qquad
 S_{E,\mathbb C}(T):=S_{E,0}(T)\otimes_{\mathbb Q}\mathbb C.
\end{equation}
Es invertible sobre $\mathbb Q((T))$ y, por tanto, sobre cada
$\mathbb Q_\ell((T))$ y $\mathbb C((T))$. Escribiremos $S_E(T)$ cuando la
base sea irrelevante. Como $K[[T]]$ es un anillo de valoración discreta para
$K=\mathbb Q,\mathbb Q_\ell$ o $\mathbb C$, su forma de Smith es
\begin{equation}\label{eq:pdf2-bsd-smith-t-adico}
 U(T)S_E(T)V(T)=\operatorname{diag}
 (T^{a_1},\ldots,T^{a_s},1,\ldots,1),
 \qquad a_i\ge1.
\end{equation}
Los exponentes $a_i$ son los mismos en todas las fibras: son las valuaciones
en $T$ de los menores no nulos de la matriz racional
\eqref{eq:pdf2-bsd-s-universal}; cambiar $\mathbb Q$ por
$\mathbb Q_\ell$ o $\mathbb C$ no cambia esas valuaciones.
Se obtiene
\begin{equation}\label{eq:pdf2-bsd-orden-smith}
 d:=\operatorname{ord}_T\det S_E(T)
 =\sum_{i=1}^sa_i
 =\sum_{q\ge1}\dim_K V_q.
\end{equation}
La filtración produce $A_q=V_q/V_{q+1}$,
$B_q=\operatorname{im}\beta_q$ e isomorfismos
\begin{equation}\label{eq:pdf2-bsd-bockstein-reducido}
 \bar\beta_q:A_q\xrightarrow{\sim}B_q,
 \qquad
 j_q:B_q\xrightarrow{\sim}A_q^\vee\otimes(I^q/I^{q+1}).
\end{equation}
Con la página regular $\bar S_E(0)$ y los signos de Koszul, todas las
páginas se pegan en
\begin{equation}\label{eq:pdf2-bsd-rboc-lt}
 R_{\rm Boc}^{\rm der}(S_E)
 =\tau_{\rm reg}\otimes\bigotimes_{q\ge1}\det(\bar\beta_q)
 =\left[T^{-d}\det S_E(T)\right]_{T=0}.
\end{equation}
Ninguna página se sustituye por el orden total $d$.

\section{Tabla completa Kato--FERS antes del determinante}

Para una historia no degenerada
$b_\lambda=[x_0\to\cdots\to x_i]$ sean
$b_{\lambda,\le j}=[x_0\to\cdots\to x_j]$ y
$b_{\lambda,\ge j}=[x_j\to\cdots\to x_i]$. El peso completo del lector
de acarreo es
\begin{equation}\label{eq:pdf2-bsd-peso-t}
 \nu_T(\gamma):=\operatorname{dep}(\gamma)
                 +\operatorname{Carr}(\gamma),
\end{equation}
de modo que $\nu_T([x_i])=0$ y toda ruta normalizada de longitud positiva
tiene $\nu_T\ge1$. La base $\mathcal B_i^K$ contiene una fila por cada
etiqueta completa $\lambda$; la base $\mathcal B_i^{\rm Iw}$ contiene el
generador $\upsilon_i(b_\lambda)$ con las mismas caras, hoja, orientación,
memoria y terminal. Por ello $\upsilon_i$ es una biyección escrita por la
misma lista de etiquetas, no elegida mediante un conteo posterior.

El mapa de cadenas es
\begin{equation}\label{eq:pdf2-bsd-phi-kato}
 \Phi_{K/{\rm FERS}}:=
 \mu_{\rm Iw}\circ(P_E^K\widehat\otimes\operatorname{car}_T)
 \circ\Delta_{\rm AW}:
 \mathcal C_E^K\longrightarrow\mathcal C_E^{\rm Iw}.
\end{equation}
Su tabla exhaustiva, parametrizada por todos los $i$, todas las etiquetas y
todos los cortes, es
\begin{equation}\label{eq:pdf2-bsd-tabla-kato-fers}
\begin{array}{c|c|c|c}
 \text{fila}&\text{corte}&\text{imagen}&\text{peso }T\\ \hline
 b_\lambda&j=i&\upsilon_i(b_\lambda)&0\\
 b_\lambda&0\le j<i&
 \epsilon_{\lambda,j}\,
 \mu_{\rm Iw}\!\left(P_E^K(b_{\lambda,\le j}),
       \operatorname{car}_T(b_{\lambda,\ge j})\right)&
 \nu_T(b_{\lambda,\ge j})\ge1\\
 \text{simplex degenerado}&\text{cualquier }j&0&--
\end{array}
\end{equation}
donde $\epsilon_{\lambda,j}$ es el signo ya fijado por
$\operatorname{or}_{\rm det}$. Equivalentemente,
\begin{equation}\label{eq:pdf2-bsd-expansion-kato-fers}
 \Phi_{K/{\rm FERS}}^i(b_\lambda)
 =\upsilon_i(b_\lambda)+
  \sum_{j=0}^{i-1}T^{\nu_T(b_{\lambda,\ge j})}
  \epsilon_{\lambda,j}\,
  \mu_{\rm Iw}\!\left(P_E^K(b_{\lambda,\le j}),
          \widehat{\operatorname{car}}_T(b_{\lambda,\ge j})\right).
\end{equation}

\begin{theorem}[Equivalencia basada Kato--FERS]
\label{thm:pdf2-bsd-equivalencia-kato-fers}
El mapa \eqref{eq:pdf2-bsd-phi-kato} es una equivalencia filtrada de
complejos perfectos, preserva la unidad y, en todas las bases completas,
\begin{equation}\label{eq:pdf2-bsd-i-plus-tb}
 [\Phi_{K/{\rm FERS}}^i](T)=I+TB_i(T),
 \qquad
 [\Phi_{K/{\rm FERS}}^i]^{-1}
 =\sum_{r\ge0}(-TB_i(T))^r.
\end{equation}
En particular,
\begin{equation}\label{eq:pdf2-bsd-rho-uno}
 \rho_E(T):=
 \prod_i\det(I+TB_i(T))^{(-1)^i},
 \qquad \rho_E(0)=1,
 \qquad
 \rho_{E,\ell}:=\rho_E\otimes_{\mathbb Q}\mathbb Q_\ell.
\end{equation}
\end{theorem}

\begin{proof}
La identidad simplicial de Alexander--Whitney, la compatibilidad de $P_E^K$
con caras y la aditividad de $\nu_T$ bajo concatenación prueban que
$\Phi_{K/{\rm FERS}}$ conmuta con los diferenciales. En cada fila de
\eqref{eq:pdf2-bsd-tabla-kato-fers}, el único término de peso cero es el
corte $j=i$; todos los demás pertenecen a $T\Lambda$. Como
$\upsilon_i$ empareja la lista completa de historias con ella misma, la
reducción módulo $T$ es la identidad en cada generador. Esto da
\eqref{eq:pdf2-bsd-i-plus-tb}; la completitud de $K[[T]]$ hace converger la
inversa. Para $i=0$ la tabla contiene sólo la unidad, de modo que la raíz
queda normalizada antes de evaluar un término principal.
\end{proof}

La equivalencia conserva por separado el ledger
\begin{equation}\label{eq:pdf2-bsd-ledger-diez}
 \mathcal J_{\rm det}=\{\deg,\operatorname{Frob},\operatorname{Ner},
 K,\operatorname{Boc},\operatorname{Sha},\operatorname{Tam},
 \operatorname{tors}^{-2},\Omega,\operatorname{exc}\}.
\end{equation}
En esas diez etiquetas se leen, respectivamente: el peso total; la norma de
Euler; el cambio basado Manin--Néron; la relación finita--singular; todas las
páginas $j_q\bar\beta_q$; el cono de Kummer; los cocientes de componentes;
la red torsional; la integral de Néron; y la orientación excepcional. La
tabla no permite cancelaciones entre etiquetas.

\section{Mapa Poitou--Tate desde Mordell--Weil a la bandera completa}

Se conserva primero la red de Mordell--Weil y sólo después se racionaliza:
\begin{equation}\label{eq:pdf2-bsd-mw-libre}
 M_{E,\mathbb Z}^{\rm MW}:=E(\mathbb Q)/E(\mathbb Q)_{\rm tors},
 \qquad
 M_E^{\rm MW}:=M_{E,\mathbb Z}^{\rm MW}\otimes_{\mathbb Z}K,
 \qquad K=\mathbb Q.
\end{equation}
$\operatorname{Surv}_{{\rm det},\mathbb Z}^{\rm MW}$ denota el submódulo
libre generado por las filas terminales integrales del sumando
Mordell--Weil.
El levantamiento de Kummer con historia completa es
\begin{equation}\label{eq:pdf2-bsd-kummer-hmt}
 \begin{aligned}
 \kappa_{E,\mathbb Z}^{\rm HMT}:M_{E,\mathbb Z}^{\rm MW}
   &\longrightarrow
   Z^1(\mathcal C_{E,\mathbb Z}^{\rm red})
   \cap\operatorname{Surv}_{{\rm det},\mathbb Z}^{\rm MW},\\
 \kappa_E^{\rm HMT}&:=
   \kappa_{E,\mathbb Z}^{\rm HMT}\otimes_{\mathbb Z}\mathbb Q:
   M_E^{\rm MW}\longrightarrow Z^1(\mathcal C_E^{\rm red}).
 \end{aligned}
\end{equation}
y $\operatorname{pg}_q$ publica la coordenada de la página $V_q$. La
extensión $\operatorname{Ext}_{\Gamma,{\rm det}}$, la multisección y la
supervivencia construyen, para cada $q$, un pegado basado
\begin{equation}\label{eq:pdf2-bsd-glue-q}
 \operatorname{gl}_q:V_q\longrightarrow
 \operatorname{Surv}_{\rm det}^{\rm MW}
\end{equation}
que conserva hoja, carry, memoria y terminal. Sus ecuaciones de
unidad--counidad son
\begin{equation}\label{eq:pdf2-bsd-unidad-counidad-paginas}
 \operatorname{pg}_rP_E^{\rm PT}\operatorname{gl}_q
 =\delta_{rq}I_{V_q},
 \qquad
 \sum_{q\ge1}\operatorname{gl}_q
     \operatorname{pg}_qP_E^{\rm PT}=I
 \quad\text{en el sumando MW superviviente}.
\end{equation}
La primera igualdad se comprueba etiqueta por etiqueta: el pegado de nivel
$q$ tiene carry $q$, cero en las otras páginas y el terminal de esa misma
historia. La segunda es la exhaustividad de
$\Sigma_{\rm det}^{\rm mult}$; la suma es finita porque la torre Bockstein
termina.

Para este submódulo póngase
\begin{equation}\label{eq:pdf2-bsd-redes-paginas-saturadas}
 V_{q,\mathbb Z}:=\operatorname{pg}_qP_E^{\rm PT}
   (\operatorname{Surv}_{{\rm det},\mathbb Z}^{\rm MW}),
 \qquad
 A_{q,\mathbb Z}:=V_{q,\mathbb Z}/V_{q+1,\mathbb Z}.
\end{equation}
En las bases terminales, $\operatorname{pg}_q$ y
$\operatorname{gl}_q$ tienen entradas $0,\pm1$ y sus dos composiciones son
las de \eqref{eq:pdf2-bsd-unidad-counidad-paginas}. Por tanto
$V_{q+1,\mathbb Z}$ es un sumando directo basado de $V_{q,\mathbb Z}$,
$A_{q,\mathbb Z}$ es libre y saturado, y
$V_q=V_{q,\mathbb Z}\otimes\mathbb Q$,
$A_q=A_{q,\mathbb Z}\otimes\mathbb Q$.
Sea $s_q:A_{q,\mathbb Z}\to V_{q,\mathbb Z}$ la sección de la base
terminal. Entonces
$V_{k,\mathbb Z}=\bigoplus_{q\geq k}s_q(A_{q,\mathbb Z})$. El desdoblamiento
de jets que conserva las $q$ ocurrencias de una fila de profundidad $q$ es
\begin{equation}\label{eq:pdf2-bsd-desdoblamiento-jets}
 \begin{aligned}
 \operatorname{Jet}_{\rm Boc,\mathbb Z}(E)
   &:=\bigoplus_{q\geq1}\bigoplus_{k=1}^{q}
       A_{q,\mathbb Z}\,e_{q,k},\\
 \mathcal U_{\rm jet}:
 \operatorname{Jet}_{\rm Boc,\mathbb Z}(E)
   &\xrightarrow{\ \sim\ }
     \bigoplus_{k\geq1}V_{k,\mathbb Z},\\
 \mathcal U_{\rm jet}((a_{q,k})_{q,k})
   &:=(v_k)_k,
   \qquad v_k:=\sum_{q\geq k}s_q(a_{q,k}).
 \end{aligned}
\end{equation}
La inversa descompone cada $v_k$ en la suma directa anterior. En particular,
\begin{equation}\label{eq:pdf2-bsd-dimension-jet}
 \operatorname{rank}_{\mathbb Z}\operatorname{Jet}_{\rm Boc,\mathbb Z}(E)
 =\sum_q q\,\operatorname{rank}_{\mathbb Z}A_{q,\mathbb Z}
 =\sum_k\operatorname{rank}_{\mathbb Z}V_{k,\mathbb Z}=d.
\end{equation}
El pegado restringe, en esas mismas bases, a
$\operatorname{gl}_{q,\mathbb Z}:V_{q,\mathbb Z}\to
\operatorname{Surv}_{{\rm det},\mathbb Z}^{\rm MW}$.

Se definen, sin usar dimensiones,
\begin{equation}\label{eq:pdf2-bsd-psi-y-lambda-q}
 \begin{aligned}
  \psi_{q,\mathbb Z}&:=
      \operatorname{pg}_qP_E^{\rm PT}\kappa_{E,\mathbb Z}^{\rm HMT}:
      M_{E,\mathbb Z}^{\rm MW}\to V_{q,\mathbb Z},\\
  \lambda_{q,\mathbb Z}&:=
      \operatorname{pr}_{\rm MW}P_E^{\rm Man}
      \operatorname{gl}_{q,\mathbb Z}:
      V_{q,\mathbb Z}\to M_{E,\mathbb Z}^{\rm MW},\\
  \psi_q&:=\operatorname{pg}_qP_E^{\rm PT}\kappa_E^{\rm HMT}:
      M_E^{\rm MW}\to V_q,\\
  \lambda_q&:=\operatorname{pr}_{\rm MW}P_E^{\rm Man}
      \operatorname{gl}_q:V_q\to M_E^{\rm MW}.
 \end{aligned}
\end{equation}
Los dominios intermedios quedan ligados por las dos compatibilidades basadas
que forman parte del pegado Poitou--Tate y que aquí se escriben sobre su
acción alcanzable:
\begin{equation}\label{eq:pdf2-bsd-kummer-glue-compatibilidad}
 \kappa_E^{\rm HMT}\operatorname{pr}_{\rm MW}P_E^{\rm Man}
       \operatorname{gl}_q=\operatorname{gl}_q,
 \qquad
 \operatorname{pr}_{\rm MW}P_E^{\rm Man}\kappa_E^{\rm HMT}
       =I_{M_E^{\rm MW}}.
\end{equation}
Las mismas dos identidades, con subíndice $\mathbb Z$, valen sobre
$M_{E,\mathbb Z}^{\rm MW}$ y $V_{q,\mathbb Z}$, porque todas las matrices
anteriores son integrales y basadas.

\begin{theorem}[Isomorfismo basado de la bandera Poitou--Tate]
\label{thm:pdf2-bsd-psi-pt}
Los mapas
\begin{equation}\label{eq:pdf2-bsd-psi-pt}
 \begin{aligned}
 \Psi_{\rm PT}:M_E^{\rm MW}&\longrightarrow
       \operatorname{Flag}_{\rm Boc}(E):=\bigoplus_{q\ge1}V_q,
 &P&\longmapsto(\psi_q(P))_{q\ge1},\\
 \Psi_{\rm PT}^{-1}:\operatorname{Flag}_{\rm Boc}(E)&\longrightarrow M_E^{\rm MW},
 &(v_q)_q&\longmapsto\sum_{q\ge1}\lambda_q(v_q)
 \end{aligned}
\end{equation}
son inversos. En consecuencia,
\begin{equation}\label{eq:pdf2-bsd-psi-pt-integral}
 \Psi_{{\rm PT},\mathbb Z}:M_{E,\mathbb Z}^{\rm MW}
 \xrightarrow{\ \sim\ }\bigoplus_{q\ge1}V_{q,\mathbb Z},
 \qquad
 \Psi_{{\rm PT},\mathbb Z}^{-1}((v_q)_q)=
 \sum_q\lambda_{q,\mathbb Z}(v_q),
\end{equation}
y, sólo después de extender escala,
\begin{equation}\label{eq:pdf2-bsd-rango-orden}
 \operatorname{rank}E(\mathbb Q)
 =\sum_{q\ge1}\dim_K V_q
 =d=\operatorname{ord}_T\det S_E(T).
\end{equation}
\end{theorem}

\begin{proof}
Aplicar \eqref{eq:pdf2-bsd-unidad-counidad-paginas} junto con
\eqref{eq:pdf2-bsd-kummer-glue-compatibilidad} da, con todos los dominios
intermedios visibles,
$\psi_r\lambda_q=\delta_{rq}I$ y
$\sum_q\lambda_q\psi_q=I_{M_E^{\rm MW}}$. Por tanto las dos fórmulas de
\eqref{eq:pdf2-bsd-psi-pt} son inversas antes de comparar dimensiones.
Las tablas terminales son integrales y basadas; el mismo cálculo, sin
denominadores, prueba \eqref{eq:pdf2-bsd-psi-pt-integral}. Así no aparece un
índice de subred al elegir después una base de Mordell--Weil.
Sólo después se toman dimensiones y se usa
\eqref{eq:pdf2-bsd-orden-smith}. El emparejamiento
$j_q\bar\beta_q$ transporta las bases saturadas a las filas de altura cuya
comparación con Néron--Tate se escribe en
\eqref{eq:pdf2-bsd-altura-nt-generadores}.
\end{proof}

\section{Localización, rango de Smith y exactitud Sha--Tamagawa}

En esta sección $v$ recorre los lugares finitos; en los de buena reducción se
tiene $\Phi_v(\mathbb F_v)=0$ y $c_v=1$, de modo que las sumas y productos
locales tienen soporte finito. El lugar infinito se reserva para la evaluación
de la diferencial de Néron en la cuarta flecha. Para cada lugar finito $v$ se
tiene el triángulo
\[
 \mathcal C_{E,v}^{\rm fin}\to\mathcal C_{E,v}\to
 \mathcal C_{E,v}^{\rm sing}\to\mathcal C_{E,v}^{\rm fin}[1]
\]
y el mapa de cadenas
\begin{equation}\label{eq:pdf2-bsd-ell-e}
 \ell_E(x,(y_v)_v)=(\operatorname{res}_vx-i_vy_v)_v.
\end{equation}
El cono reducido es
\begin{equation}\label{eq:pdf2-bsd-cono-reducido}
 \mathcal C_E^{\rm loc,red}:=
 q_{\rm root}\operatorname{Cone}(\ell_E)[-1].
\end{equation}
En grados de cociclos, un elemento se escribe $(x,(y_v),(h_v))$ y
\begin{equation}\label{eq:pdf2-bsd-diferencial-cono}
 d(x,(y_v),(h_v))=
 (dx,(dy_v),(\operatorname{res}_vx-i_vy_v-dh_v)_v).
\end{equation}

Tras separar por \eqref{eq:pdf2-bsd-psi-pt} la red libre de Mordell--Weil
y su dual, los bloques FS están contraídos por
\eqref{eq:pdf2-bsd-hfs-global}. En el retículo restante, el inverso
\eqref{eq:pdf2-bsd-j-e-inverso-racional} y sus dos comprobaciones
\eqref{eq:pdf2-bsd-j-e-dos-lados} dan, sin inferir aciclicidad sólo desde la
dualidad,
\begin{equation}\label{eq:pdf2-bsd-rango-completo}
 D_E\otimes\mathbb Q:M_E\otimes\mathbb Q
 \xrightarrow{\sim}N_E\otimes\mathbb Q.
\end{equation}
La forma normal integral da
\begin{equation}\label{eq:pdf2-bsd-smith-finito}
 UD_EV=\operatorname{diag}(s_1,\ldots,s_t),
 \qquad
 F_E:=\operatorname{coker}D_E,
 \qquad |F_E|=\prod_{i=1}^t|s_i|<\infty.
\end{equation}

Sea
\begin{equation}\label{eq:pdf2-bsd-pi-componentes}
 \pi_\Phi:N_E\longrightarrow\bigoplus_v\Phi_v(\mathbb F_v)
\end{equation}
la reducción de cada coordenada local módulo
$\mathcal C_{E,v}^{\rm fin}+d\mathcal C_{E,v}$; se tiene
$\pi_\Phi D_E=0$. Para una clase
$[\xi]\in\operatorname{Sha}(E)$, elíjanse sus primitivas finitas
$y_v$ y homotopías $h_v$ con
$\operatorname{res}_v\xi-i_vy_v=dh_v$. La fórmula
\begin{equation}\label{eq:pdf2-bsd-iota-sha-cociclos}
 \begin{aligned}
 n_\xi&:=\operatorname{pr}_{N_E}\operatorname{pr}_{\mathcal L_E}
       \Theta_E(\xi,(y_v),(h_v)),\\
 \iota_{\rm Sha}[\xi]&:=[n_\xi]\in F_E
 \end{aligned}
\end{equation}
es independiente de las elecciones por
\eqref{eq:pdf2-bsd-diferencial-cono}.

Para hacer material la sobreyectividad final, sea
$\widetilde\phi_v$ el representante terminal basado de
$\phi_v\in\Phi_v(\mathbb F_v)$. La suma de counidades se corrige en la
unidad común y se levanta por el operador estructural:
\begin{equation}\label{eq:pdf2-bsd-lift-componentes}
 \operatorname{Lift}_\Phi((\phi_v)_v):=
 \operatorname{pr}_{N_E}\operatorname{Ext}_{\Gamma,{\rm det}}
 \left((\widetilde\phi_v)_v,
 -\widehat\Pi_0\!\left(\sum_v
       \varepsilon_{\rm AW}(\widetilde\phi_v)\right)\right).
\end{equation}
La relación de corestricción de
$\operatorname{Rel}_{\rm det}$ hace de esta expresión un cociclo y da
\begin{equation}\label{eq:pdf2-bsd-lift-dos-identidades}
 \pi_\Phi\operatorname{Lift}_\Phi=I,
 \qquad
 \operatorname{im}D_E=
 \ker\pi_\Phi\cap\ker(\text{clase global}).
\end{equation}

\begin{theorem}[Sucesión exacta finita Sha--Tamagawa]
\label{thm:pdf2-bsd-sha-tamagawa}
Los mapas de cociclos anteriores inducen
\begin{equation}\label{eq:pdf2-bsd-sha-tamagawa}
 0\longrightarrow\operatorname{Sha}(E)
 \xrightarrow{\ \iota_{\rm Sha}\ }F_E
 \xrightarrow{\ \partial_{\rm loc}\ }
 \bigoplus_v\Phi_v(\mathbb F_v)
 \longrightarrow0,
 \qquad
 \partial_{\rm loc}[n]:=\pi_\Phi(n).
\end{equation}
En particular, $\operatorname{Sha}(E)$ es finito y
\begin{equation}\label{eq:pdf2-bsd-cardinal-sts}
 |F_E|=|\operatorname{Sha}(E)|\prod_vc_v,
 \qquad c_v:=|\Phi_v(\mathbb F_v)|.
\end{equation}
\end{theorem}

\begin{proof}
La igualdad $\pi_\Phi D_E=0$ hace bien definida a
$\partial_{\rm loc}$. Si \eqref{eq:pdf2-bsd-iota-sha-cociclos} es una
frontera, su componente global es una frontera; así $\iota_{\rm Sha}$ es
inyectiva. Una clase de $F_E$ está en $\ker\partial_{\rm loc}$ si y sólo
si todas sus coordenadas singulares admiten representantes finitos; por
\eqref{eq:pdf2-bsd-diferencial-cono}, esto equivale a estar en la imagen de
$\iota_{\rm Sha}$. Finalmente,
\eqref{eq:pdf2-bsd-lift-dos-identidades} da una preimagen de cada tupla de
componentes y prueba la sobreyectividad. Los términos adyacentes libres de
la sucesión larga son precisamente los dos sumandos separados por
$\Psi_{\rm PT}$ y $X^{\rm orth}$; los términos adyacentes FS son
contractibles por $H_{\rm FS}$. No se ha supuesto su anulación. La finitud
de $F_E$ procede de \eqref{eq:pdf2-bsd-rango-completo}; la inyección da la
 finitud de $\operatorname{Sha}(E)$ y tomar órdenes da
 \eqref{eq:pdf2-bsd-cardinal-sts}.
\end{proof}

\begin{corollary}[Descarga primaria completa y recomposición integral]
\label{cor:pdf2-bsd-sha-tamagawa-todos-ell}
Para cada primo $\ell$, el cambio de base de los mapas de cociclos produce
la sucesión exacta
\begin{equation}\label{eq:pdf2-bsd-sha-tamagawa-ell}
 0\longrightarrow\operatorname{Sha}(E)[\ell^\infty]
 \xrightarrow{\ \iota_{{\rm Sha},\ell}\ }F_{E,\ell}
 \xrightarrow{\ \partial_{{\rm loc},\ell}\ }
 \bigoplus_v\Phi_v(\mathbb F_v)[\ell^\infty]
 \longrightarrow0,
\end{equation}
donde
\begin{equation}\label{eq:pdf2-bsd-mapas-ell-generadores}
 \iota_{{\rm Sha},\ell}([\xi]\otimes1)
   =[n_\xi\otimes1],
 \qquad
 \partial_{{\rm loc},\ell}([n\otimes1])
   =(\pi_{\Phi,v}(n)\otimes1)_v.
\end{equation}
Además, la suma de todas las sucesiones
\eqref{eq:pdf2-bsd-sha-tamagawa-ell} es exactamente
\eqref{eq:pdf2-bsd-sha-tamagawa}; no es sólo su parte $3$-primaria.
\end{corollary}

\begin{proof}
La matriz $D_E$, el cociclo $n_\xi$, la reducción $\pi_\Phi$ y
$\operatorname{Lift}_\Phi$ son integrales por
\eqref{eq:pdf2-bsd-complejos-base-change}--
\eqref{eq:pdf2-bsd-theta-base-change}. Como $\mathbb Z_\ell$ es plana sobre
$\mathbb Z$, tensorizar la sucesión exacta finita
\eqref{eq:pdf2-bsd-sha-tamagawa} conserva exactitud. Para un grupo finito
$A$ se tiene canónicamente
$A\otimes\mathbb Z_\ell=A[\ell^\infty]$; esto identifica los tres términos
y da las fórmulas sobre generadores
\eqref{eq:pdf2-bsd-mapas-ell-generadores}. La forma de Smith
\eqref{eq:pdf2-bsd-smith-finito} escribe
$s_i=\prod_\ell\ell^{v_\ell(s_i)}$, de modo que
\[
 F_E\cong\bigoplus_\ell F_{E,\ell},\qquad
 \operatorname{Sha}(E)\cong
   \bigoplus_\ell\operatorname{Sha}(E)[\ell^\infty],\qquad
 \Phi_v(\mathbb F_v)\cong
   \bigoplus_\ell\Phi_v(\mathbb F_v)[\ell^\infty].
\]
Sólo intervienen finitos primos y finitos lugares porque todos los grupos
son cocientes del grupo finito $F_E$. Sumar las sucesiones primarias recupera
la integral. Tomando órdenes primero primo a primo y luego multiplicando,
\begin{equation}\label{eq:pdf2-bsd-cardinal-sts-todos-ell}
 \prod_\ell|F_{E,\ell}|
 =\prod_\ell|\operatorname{Sha}(E)[\ell^\infty]|
   \prod_v\prod_\ell|\Phi_v(\mathbb F_v)[\ell^\infty]|
 =|\operatorname{Sha}(E)|\prod_vc_v.
\end{equation}
Así quedan descargados materialmente todos los coeficientes y todos los
factores de Tamagawa.
\end{proof}

\section{Comparador analítico--aritmético basado}

La coordenada analítica no se postula desde el orden buscado. Sea
$\lambda_3=\log4$, correspondiente al generador ciclotómico basado
$\gamma_3=4$, y defínase en un disco $U_E$ centrado en $1$
\begin{equation}\label{eq:pdf2-bsd-t-e-construido}
 T_E(s):=\frac{\exp(\lambda_3(s-1))-1}{\lambda_3}
 =(s-1)+\frac{\lambda_3}{2}(s-1)^2+O((s-1)^3).
\end{equation}
Por cálculo directo,
\begin{equation}\label{eq:pdf2-bsd-t-e-normalizacion-probada}
 T_E(1)=0,
 \qquad T_E'(s)=\exp(\lambda_3(s-1)),
 \qquad T_E'(1)=1.
\end{equation}
El teorema de la función inversa proporciona, después de reducir $U_E$ si
fuera necesario, una carta biholomorfa $s\leftrightarrow T_E(s)$.
Para comparar todas las fibras, tómese
$\gamma_\ell=1+\ell$ si $\ell$ es impar y $\gamma_2=5$, y póngase
\begin{equation}\label{eq:pdf2-bsd-t-ell-construido}
 T_{E,\ell}(s):=
 \frac{\exp(\log(\gamma_\ell)(s-1))-1}{\log(\gamma_\ell)}.
\end{equation}
Las cartas se relacionan por un único automorfismo formal
$T_{E,\ell}=\vartheta_{\ell,3}(T_E)$ con
$\vartheta_{\ell,3}(0)=0$ y $\vartheta_{\ell,3}'(0)=1$. Por tanto ni el
orden en $T$ ni el coeficiente principal cambian al pasar de $3$ a $\ell$.

La continuación analítica que llega a esta carta tampoco se declara. Por el
teorema de modularidad, sea
\[
 f_E(z)=\sum_{m\geq1}a_m(E)e^{2\pi imz}
\]
la forma nueva normalizada de peso $2$ y nivel $N_E$ asociada a $E$, y sea
$\varepsilon_E\in\{\pm1\}$ el signo fijado por
\begin{equation}\label{eq:pdf2-bsd-transformacion-modular}
 g_E(y):=f_E\!\left(\frac{iy}{\sqrt{N_E}}\right),
 \qquad
 g_E(1/y)=\varepsilon_Ey^2g_E(y).
\end{equation}
Para $\Re s>3/2$, integración término a término da
\begin{equation}\label{eq:pdf2-bsd-mellin-semiplano}
 \Lambda_E(s):=N_E^{s/2}(2\pi)^{-s}\Gamma(s)L(E,s)
 =\int_0^\infty g_E(y)y^s\,\frac{dy}{y}.
\end{equation}
Se divide la integral en $(0,1)$ y $(1,\infty)$ y, en la primera parte, se
efectúa el cambio $y=1/t$. La segunda igualdad de
\eqref{eq:pdf2-bsd-transformacion-modular} produce la fórmula material
\begin{equation}\label{eq:pdf2-bsd-mellin-continuada}
 \Lambda_E^{\rm cont}(s)=
 \int_1^\infty g_E(y)
       \bigl(y^s+\varepsilon_Ey^{2-s}\bigr)\,\frac{dy}{y}.
\end{equation}
La función $g_E(y)$ decrece exponencialmente cuando $y\to\infty$; por ello la
integral de \eqref{eq:pdf2-bsd-mellin-continuada}, y todas sus derivadas en
$s$, convergen uniformemente en compactos. Así define una función entera y
coincide con \eqref{eq:pdf2-bsd-mellin-semiplano} donde la serie de Dirichlet
converge. Puesto que
$N_E^{s/2}(2\pi)^{-s}\Gamma(s)$ es holomorfa y no nula cerca de $s=1$, la
división por ese factor construye, sin usar su orden de anulación, el germen
holomorfo de $L(E,s)$ en un disco $U_E$ centrado en $1$.

Para evitar mezclar una fibra $\ell$-ádica con una matriz compleja, se toma
primero el polinomio entero, independiente de $\ell$,
\begin{equation}\label{eq:pdf2-bsd-polinomio-local-entero}
 P_{E,p}(X):=\det(I-X\operatorname{Frob}_p\mid V_\ell E^{I_p})
 \in\mathbb Z[X].
\end{equation}
Para $p\nmid N_E$, sean $\alpha_p,\beta_p$ las raíces de $P_{E,p}$. Se toma
$W_p=\mathbb C^2$ con base ortonormal y la realización normal
\begin{equation}\label{eq:pdf2-bsd-realizacion-normal-frobenius}
 F_p^{\mathbb C}:=\operatorname{diag}(\alpha_p,\beta_p),
 \qquad |\alpha_p|=|\beta_p|=p^{1/2},
 \qquad \|F_p^{\mathbb C}\|_1=2p^{1/2}.
\end{equation}
Para los primos malos se elige de igual modo una realización compleja del
polinomio local; son finitos y se conservan como bloques separados. No se usa
la matriz compañera, cuya norma no estaría controlada uniformemente por el
módulo de sus autovalores.

Póngase $\mathscr H_E:=\widehat\bigoplus_{p\nmid N_E}W_p$ y, para $X>0$,
\begin{equation}\label{eq:pdf2-bsd-truncados-fredholm}
 K_{E,X}(s):=\bigoplus_{\substack{p\leq X\\p\nmid N_E}}
          p^{-s}F_p^{\mathbb C},
 \qquad
 L_X(E,s):=\det(I-K_{E,X}(s))^{-1}
       \prod_{p\mid N_E}P_{E,p}(p^{-s})^{-1}.
\end{equation}
Si $K_0$ es un compacto de $\{\Re s>3/2\}$ y
$\sigma_0:=\inf_{s\in K_0}\Re s$, entonces
\begin{equation}\label{eq:pdf2-bsd-cota-traza}
 \sup_{s\in K_0}\|K_{E,X'}(s)-K_{E,X}(s)\|_1
 \leq C_E\sum_{X<p\leq X'}p^{1/2-\sigma_0}
 \xrightarrow[X,X'\to\infty]{}0.
\end{equation}
Así $K_{E,X}\to K_E$ en norma traza, localmente uniformemente, y la
continuidad del determinante de Fredholm prueba, no define,
\begin{equation}\label{eq:pdf2-bsd-fredholm-euler}
 \det_F(I-K_E(s))^{-1}
 \prod_{p\mid N_E}P_{E,p}(p^{-s})^{-1}
 =\lim_{X\to\infty}L_X(E,s)=L(E,s).
\end{equation}

Para pasar del producto infinito a un complejo finito se conserva primero el
operador, no sólo su determinante. Sean $\mathscr H_E^+$ y
$\mathscr H_E^-$ dos copias basadas de
$\mathscr H_E\oplus\bigoplus_{p\mid N_E}W_p$. En el semiplano
$\Re s>3/2$ se define
\begin{equation}\label{eq:pdf2-bsd-operadores-euler-inversos}
 \begin{aligned}
  \mathscr A_{E,X}(s)&:=
   (I-K_{E,X}(s))^{-1}\oplus
   \bigoplus_{p\mid N_E}(I-p^{-s}F_p^{\mathbb C})^{-1},\\
  \mathscr A_E(s)&:=\lim_{X\to\infty}\mathscr A_{E,X}(s):
       \mathscr H_E^+\longrightarrow\mathscr H_E^-.
 \end{aligned}
\end{equation}
La identidad del resolvente y \eqref{eq:pdf2-bsd-cota-traza} dan, para cada
compacto $K_0$ del semiplano,
\begin{equation}\label{eq:pdf2-bsd-resolvente-convergencia-traza}
 \sup_{s\in K_0}
 \|\mathscr A_{E,X'}(s)-\mathscr A_{E,X}(s)\|_1\longrightarrow0.
\end{equation}
Sea $\mathcal U_{\rm bas}:\mathscr H_E^+\to\mathscr H_E^-$ la identificación unitaria que
envía cada vector de la base positiva al vector homónimo de la base
negativa. Los determinantes siguientes significan los de
$\mathcal U_{\rm bas}^{-1}\mathscr A_{E,X}=I+\mathcal L^1$ y
$\mathcal U_{\rm bas}^{-1}\mathscr A_E=I+\mathcal L^1$. Entonces
\begin{equation}\label{eq:pdf2-bsd-det-a-x-l-x}
 \det_F(\mathcal U_{\rm bas}^{-1}\mathscr A_{E,X}(s))=L_X(E,s),
 \qquad
 \det_F(\mathcal U_{\rm bas}^{-1}\mathscr A_E(s))=L(E,s).
\end{equation}
La segunda igualdad es precisamente el límite de
\eqref{eq:pdf2-bsd-fredholm-euler}; el signo es el correcto porque se toma el
inverso de $I-K_E$.

La corestricción finita se construye a partir de los terminales, sin publicar
infinitos primos en una matriz finita. Sean $\mathcal I_E^+$ y
$\mathcal I_E^-$ las listas terminales de las bases integrales de $M_E$ y
$N_E$, respectivamente; son finitas por
\eqref{eq:pdf2-bsd-sumando-fs-y-lattice}. Fijado $\sigma_*>3/2$, el
entrelazador local--global las inserta en el espacio de Euler por
\begin{equation}\label{eq:pdf2-bsd-embedding-terminal-euler}
 \iota_{\rm Eul}^\pm(e_\eta):=
 \left(p^{-\sigma_*}\operatorname{res}_p^\pm
       \operatorname{Surv}_E(e_\eta)\right)_p
 \oplus\left(\operatorname{res}_p^\pm
       \operatorname{Surv}_E(e_\eta)\right)_{p\mid N_E}.
\end{equation}
La suma buena es cuadrado--sumable por la misma cota que
\eqref{eq:pdf2-bsd-cota-traza}. Se ponen
$\tau_\eta^\pm:=\iota_{\rm Eul}^\pm(e_\eta)$; la independencia de las filas
terminales da las counidades duales $\{\epsilon_{\eta}^{\pm}\}$, de modo que
$\epsilon_\eta^\pm(\tau_{\eta'}^\pm)=\delta_{\eta\eta'}$. Las sumas
\begin{equation}\label{eq:pdf2-bsd-proyecciones-terminales-finitas}
 P_E^\pm:=\sum_{\eta\in\mathcal I_E^\pm}
      \tau_\eta^\pm\epsilon_\eta^\pm,
 \qquad Q_E^\pm:=I-P_E^\pm
\end{equation}
son proyecciones basadas de rango finito; sus índices se fijan por las
etiquetas terminales y no por el rango de Mordell--Weil ni por un orden de
anulación.

La coordenada $\deg$ del ledger descompone
\begin{equation}\label{eq:pdf2-bsd-graduacion-fredholm}
 \mathscr H_E^\pm=\widehat\bigoplus_i\mathscr H_E^{\pm,i},
 \qquad
 \mathcal I_E^{\pm,i}:=\{\eta\in\mathcal I_E^\pm:\deg\eta=i\},
 \qquad
 P_E^{\pm,i}:=
 \sum_{\eta\in\mathcal I_E^{\pm,i}}\tau_\eta^\pm\epsilon_\eta^\pm.
\end{equation}
Las relaciones locales conservan $\deg$ salvo el diferencial $i\mapsto i+1$;
por tanto
$\mathscr A_E=\bigoplus_i\mathscr A_E^i$ con
$\mathscr A_E^i:\mathscr H_E^{+,i}\to\mathscr H_E^{-,i+1}$. Los cuatro
bloques de \eqref{eq:pdf2-bsd-bloques-fredholm}, su Schur y sus truncaciones
se toman grado a grado, y por definición
\begin{equation}\label{eq:pdf2-bsd-schur-graduado}
 \mathscr D_E=\bigoplus_i\mathscr D_E^i,
 \qquad
 \mathscr D_E^i:=A_{00}^i-A_{01}^i(A_{11}^i)^{-1}A_{10}^i.
\end{equation}

La continuación Fredholm se escribe antes de usarla. En la carta de Euler
$U_0$ se pone $\mathscr A_{E,0}:=\mathscr A_E$. Las cartas conservan un
complemento basado común $\mathscr Q_E^\pm$ y cambian sólo la familia finita
de terminales. Si
$\tau_{\eta,\alpha}^\pm:=
\operatorname{Ext}_{\Gamma,{\rm det},\alpha}(\tau_\eta^\pm)$ y
$\epsilon_{\eta,\alpha}^\pm$ es su base dual, se define
\begin{equation}\label{eq:pdf2-bsd-cambios-fredholm-cartas}
 C_{\alpha\beta}^\pm:=
 I_{\mathscr Q_E^\pm}\oplus
 \sum_{\eta\in\mathcal I_E^\pm}
 \tau_{\eta,\beta}^\pm\epsilon_{\eta,\alpha}^\pm,
 \qquad C_{\alpha\beta}^\pm-I\in\mathcal L^1,
 \qquad
 C_{\beta\gamma}^\pm C_{\alpha\beta}^\pm=C_{\alpha\gamma}^\pm.
\end{equation}
La diferencia respecto de la identidad tiene rango a lo sumo
$|\mathcal I_E^\pm|$ y es, por tanto, de clase traza; la holomorfía se sigue
de que cada una de las finitas coordenadas terminales de
$\operatorname{Ext}_{\Gamma,{\rm det},\alpha}$ es holomorfa. La última
igualdad se verifica en $\mathscr Q_E^\pm$ y en cada
$\tau_{\eta,\alpha}^\pm$. Se define
recursivamente en los solapamientos
\begin{equation}\label{eq:pdf2-bsd-continuacion-fredholm-cartas}
 \mathscr A_{E,\beta}:=
 C_{\alpha\beta}^-\mathscr A_{E,\alpha}
 (C_{\alpha\beta}^+)^{-1}.
\end{equation}
Los operadores $C_{\alpha\beta}^\pm$ son holomorfos porque cada coeficiente
es una coordenada de $\operatorname{Ext}_{\Gamma,{\rm det}}$; el cociclo de
\eqref{eq:pdf2-bsd-cambios-fredholm-cartas} prueba que
\eqref{eq:pdf2-bsd-continuacion-fredholm-cartas} no depende de la cadena de
cartas. Ésta es la familia determinante--clase que se denota en adelante por
$\mathscr A_E(s)$. Con
$C_{0\alpha}^\pm:=C_{\alpha-1,\alpha}^\pm\cdots C_{0,1}^\pm$ se transportan
además
$P_{E,\alpha}^{\pm,i}:=
C_{0\alpha}^\pm P_E^{\pm,i}(C_{0\alpha}^\pm)^{-1}$ y
$\mathscr H_{E,\alpha}^{\pm,i}:=
C_{0\alpha}^\pm\mathscr H_E^{\pm,i}$; de este modo todos los Schur
$\mathscr D_{E,\alpha}^i$ son los de
\eqref{eq:pdf2-bsd-schur-graduado}, no símbolos nuevos. Tras reducir cada
carta, el bloque
\begin{equation}\label{eq:pdf2-bsd-bloques-fredholm}
 A_{ab}(s):=P_a^-\mathscr A_E(s)P_b^+,
 \qquad P_0^\pm:=P_E^\pm,\quad P_1^\pm:=Q_E^\pm,
 \qquad A_{11}(s):Q_E^+\mathscr H_E^+\xrightarrow{\sim}
                       Q_E^-\mathscr H_E^-
\end{equation}
es invertible. Ésta es la reducción finita estándar del operador Fredholm;
la elección concreta de $P_E^\pm$ viene dada por
\eqref{eq:pdf2-bsd-proyecciones-terminales-finitas}. La base fija también
la identificación unitaria
$U_{11}:Q_E^+\mathscr H_E^+\to Q_E^-\mathscr H_E^-$; por construcción
$U_{11}^{-1}A_{11}=I+\mathcal L^1$, y sólo a este operador se aplica
$\det_F$. Las copias $+$ y $-$ ya son la totalización par--impar; no se
aplica un segundo signo alternado al determinante del operador inverso.

Para cada truncación que contiene el soporte de esos terminales se definen
materialmente
\begin{equation}\label{eq:pdf2-bsd-cor-lift-truncados}
 \begin{aligned}
  \operatorname{Cor}_{E,X}&:=P_E^-,\\
  \operatorname{Lift}_{E,X}(v)&:=
       v-A_{11,X}^{-1}A_{10,X}v,\\
  \mathscr D_{E,X}(s)&:=
       \operatorname{Cor}_{E,X}\mathscr A_{E,X}(s)
       \operatorname{Lift}_{E,X}
       =A_{00,X}-A_{01,X}A_{11,X}^{-1}A_{10,X}.
 \end{aligned}
\end{equation}
Por multiplicación de bloques,
\begin{equation}\label{eq:pdf2-bsd-eliminacion-schur}
 \begin{pmatrix}I&-A_{01,X}A_{11,X}^{-1}\\0&I\end{pmatrix}
 \mathscr A_{E,X}
 \begin{pmatrix}I&0\\-A_{11,X}^{-1}A_{10,X}&I\end{pmatrix}
 =\begin{pmatrix}\mathscr D_{E,X}&0\\0&A_{11,X}\end{pmatrix}.
\end{equation}
Así la corestricción no es nominal: su lift, su inversa sobre el complemento
y su diferencial finito están escritos. La convergencia del resolvente
implica
\begin{equation}\label{eq:pdf2-bsd-corestriccion-limite}
 \mathscr D_E(s):=\lim_{X\to\infty}\mathscr D_{E,X}(s)
 =A_{00}-A_{01}A_{11}^{-1}A_{10},
 \qquad
 \|\mathscr D_{E,X}-\mathscr D_E\|_1\to0.
\end{equation}
El determinante del complemento no se descarta: por
\eqref{eq:pdf2-bsd-eliminacion-schur},
\begin{equation}\label{eq:pdf2-bsd-determinante-schur-complemento}
 \det_F(\mathcal U_{\rm bas}^{-1}\mathscr A_E)
 =\det_F(U_{11}^{-1}A_{11})\det(\mathscr D_E).
\end{equation}

Sea $\mathcal O(U_E)$ el anillo de funciones holomorfas del disco central.
Los espacios racionales basados $V_0^i,W_0^{i+1}$ de
\eqref{eq:pdf2-bsd-s-universal} proporcionan
\begin{equation}\label{eq:pdf2-bsd-dominios-analiticos}
 \mathscr V_E^i:=V_0^i\otimes_{\mathbb Q}\mathcal O(U_E),
 \qquad
 \mathscr W_E^{i+1}:=W_0^{i+1}\otimes_{\mathbb Q}\mathcal O(U_E).
\end{equation}
Las filas de $\operatorname{Rel}_{\rm det}$ construyen, sobre cada terminal
$\eta$ y en cada grado, los vectores del modelo común
$v_{\eta,i}:=\upsilon_i(b_\eta)$,
$w_{\eta,i+1}:=\upsilon_{i+1}(b_\eta^{\rm term})$ y los isomorfismos basados
\begin{equation}\label{eq:pdf2-bsd-jvw-generadores}
 \begin{aligned}
  J_{E,V,i}(\tau_{\eta,i}^+)&:=v_{\eta,i},\\
  J_{E,W,i+1}(\tau_{\eta,i+1}^-)&:=w_{\eta,i+1}.
 \end{aligned}
\end{equation}
La relación diferencial de esa misma fila, aplicada a cada generador, es
\begin{equation}\label{eq:pdf2-bsd-entrelazamiento-analitico-smith}
 \boxed{
 J_{E,W,i+1}(s)\,\mathscr D_E^i(s)
 =S_E^i(T_E(s))\,J_{E,V,i}(s).}
\end{equation}
Éste es el entrelazador efectivo entre el Schur Fredholm corestringido y la
matriz de Smith; no es una igualdad de determinantes postulada.

Para tipar la continuación, sea $\mathfrak U_E$ un dominio conexo que contiene
un abierto del semiplano $\Re s>3/2$ y el disco $U_E$, y elíjase una cadena
finita de cartas $\{U_\alpha\}_{\alpha=0}^A$ de
$\operatorname{Ext}_{\Gamma,{\rm det}}$, con $U_A=U_E$. En $U_\alpha$ se
ponen
$\mathscr P_{E,\alpha}^i:=
P_{E,\alpha}^{+,i}\mathscr H_{E,\alpha}^{+,i}$ y
$\mathscr D_{E,\alpha}^i$ igual al Schur de
\eqref{eq:pdf2-bsd-corestriccion-limite}. Los mapas al modelo común se
definen, sin elección adicional, por
\begin{equation}\label{eq:pdf2-bsd-j-cartas-modelo-comun}
 J_{E,\alpha,i}(\tau_{\eta,\alpha,i}):=v_{\eta,i}
 \quad\text{en el dominio},
 \qquad
 J_{E,\alpha,i+1}(\tau_{\eta,\alpha,i+1}):=w_{\eta,i+1}
 \quad\text{en el codominio}.
\end{equation}
Son holomorfos porque las bases de
\eqref{eq:pdf2-bsd-cambios-fredholm-cartas} lo son. En un solapamiento se
define, siempre en el mismo grado,
\begin{equation}\label{eq:pdf2-bsd-transiciones-atlas}
 G_{\alpha\beta,i}:=J_{E,\beta,i}^{-1}J_{E,\alpha,i}:
 \mathscr P_{E,\alpha}^i\big|_{U_\alpha\cap U_\beta}
 \longrightarrow
 \mathscr P_{E,\beta}^i\big|_{U_\alpha\cap U_\beta}.
\end{equation}
La composición literal de estas matrices prueba
\begin{equation}\label{eq:pdf2-bsd-cociclo-atlas}
 G_{\beta\gamma,i}G_{\alpha\beta,i}=G_{\alpha\gamma,i},
 \qquad G_{\alpha\alpha,i}=I,
 \qquad G_{\beta\alpha,i}=G_{\alpha\beta,i}^{-1},
\end{equation}
y el entrelazamiento en las dos cartas da la compatibilidad de cadena
\begin{equation}\label{eq:pdf2-bsd-cociclo-cadena}
 G_{\alpha\beta,i+1}\mathscr D_{E,\alpha}^i
 =\mathscr D_{E,\beta}^iG_{\alpha\beta,i}.
\end{equation}
Por tanto
$g_{\alpha\beta}:=\prod_i\det(G_{\alpha\beta,i})^{(-1)^i}$ satisface el
cociclo triple. Si
\begin{equation}\label{eq:pdf2-bsd-tau-atlas}
 \tau_E^{\rm atlas}(s):=
 \prod_{\alpha=0}^{A-1}g_{\alpha,\alpha+1}(s),
\end{equation}
entonces $\tau_E^{\rm atlas}$ es exactamente el transporte determinantal
desde la carta de Euler a la carta central.

Se denotan por $J_{E,V}$ y $J_{E,W}$ las totalizaciones par--impar, con
orden de Koszul, de los mapas de
\eqref{eq:pdf2-bsd-entrelazamiento-analitico-smith}; sus determinantes son
los productos alternados de los determinantes por grado. En la carta central
se define, y no antes, la unidad conjunta
\begin{equation}\label{eq:pdf2-bsd-u-e-construida}
 u_E(s):=\tau_E^{\rm atlas}(s)
 \det_F(U_{11}^{-1}A_{11}(s))
 \frac{\det J_{E,V}(s)}
      {\det J_{E,W}(s)}.
\end{equation}
Todos sus factores son holomorfos e invertibles al reducir $U_E$. La
normalización no se impone factor a factor. Sea $\mathbf e_{\rm cent}$ el
generador de la línea determinante del modelo común. Al efectuar las dos
operaciones triangulares de \eqref{eq:pdf2-bsd-eliminacion-schur}, transportar
por \eqref{eq:pdf2-bsd-cambios-fredholm-cartas} y aplicar
\eqref{eq:pdf2-bsd-j-cartas-modelo-comun}, la acción conjunta en $s=1$ es,
generador por generador,
\begin{equation}\label{eq:pdf2-bsd-u-e-normalizacion-probada}
 \begin{aligned}
 &\tau_E^{\rm atlas}(1)
 \det_F(U_{11}^{-1}A_{11}(1))
 \frac{\det J_{E,V}(1)}
      {\det J_{E,W}(1)}\,\mathbf e_{\rm cent}\\
 &\qquad=
 \bigotimes_{i,\eta}
 \bigl(\epsilon_{\eta,i}^-(\tau_{\eta,i}^-)
       \epsilon_{\eta,i}^+(\tau_{\eta,i}^+)^{-1}\bigr)^{(-1)^i}
 \mathbf e_{\rm cent}
 =\mathbf e_{\rm cent},
 \qquad \boxed{u_E(1)=1}.
 \end{aligned}
\end{equation}
La primera igualdad es la identidad de conservación de la base común del
comparador de $\mathfrak G_{\rm det}(E)$, ahora expandida por sus filas; la
segunda usa
$\epsilon_{\eta,i}^\pm(\tau_{\eta,i}^\pm)=1$. No se ha normalizado ninguno
de los cuatro factores por separado ni se ha usado el valor buscado.

\begin{proposition}[Procedencia FERS de la factorización analítica]
\label{prop:pdf2-bsd-fers-factorizacion-owner}
En el disco $U_E$ se tiene la identidad de secciones determinantales
\begin{equation}\label{eq:pdf2-bsd-factorizacion-analitica}
 \boxed{L(E,s)=u_E(s)\det S_E(T_E(s)).}
\end{equation}
La igualdad es la publicación del lector FERS de
$\mathfrak G_{\rm det}(E)$ y no una igualdad postulada después de
calcular el orden o el término principal.
\end{proposition}

\begin{proof}
En el semiplano absoluto,
\eqref{eq:pdf2-bsd-det-a-x-l-x} identifica el límite Fredholm con $L(E,s)$.
La eliminación literal \eqref{eq:pdf2-bsd-eliminacion-schur} y la continuidad
en norma traza dan
$L(E,s)=\det_F(U_{11}^{-1}A_{11}(s))\det\mathscr D_E(s)$.
Tomar cuñas ordenadas en
\eqref{eq:pdf2-bsd-entrelazamiento-analitico-smith} da
\[
 \det\mathscr D_E(s)
 =\frac{\det J_{E,V}(s)}
        {\det J_{E,W}(s)}
   \det S_E(T_E(s)).
\]
Las ecuaciones \eqref{eq:pdf2-bsd-cociclo-atlas} y
\eqref{eq:pdf2-bsd-cociclo-cadena} transportan esta sección por las cartas;
el producto exacto de sus transiciones es $\tau_E^{\rm atlas}$. Por tanto, en
$U_A=U_E$, el coeficiente que queda es precisamente
\eqref{eq:pdf2-bsd-u-e-construida}, lo que prueba
\eqref{eq:pdf2-bsd-factorizacion-analitica}. La unicidad de la continuación
Mellin de \eqref{eq:pdf2-bsd-mellin-continuada} excluye otra sección. Por
último, \eqref{eq:pdf2-bsd-u-e-normalizacion-probada} prueba conjuntamente
$u_E(1)=1$. Ningún paso usa $d$, el rango, $\operatorname{Sha}(E)$ ni el
término principal.
\end{proof}

Defínase, a partir de la forma de Smith y no del orden analítico,
\begin{equation}\label{eq:pdf2-bsd-a-e-smith}
 a_E(T):=T^{-d}\det S_E(T).
\end{equation}
Por \eqref{eq:pdf2-bsd-rboc-lt}, $a_E$ es holomorfa en $T=0$ y
\begin{equation}\label{eq:pdf2-bsd-a-e-cero}
 a_E(0)=R_{\rm Boc}^{\rm der}(S_E)\ne0.
\end{equation}
Sustituyendo esta identidad y
\eqref{eq:pdf2-bsd-t-e-construido} en la factorización analítica se obtiene
materialmente
\begin{equation}\label{eq:pdf2-bsd-expansion-local-deducida}
 \begin{aligned}
 L(E,s)
  &=(s-1)^d\,U_E^{\rm lead}(s),\\
 U_E^{\rm lead}(s)
  &:=u_E(s)
     \left(\frac{T_E(s)}{s-1}\right)^d
     a_E(T_E(s)),\\
 U_E^{\rm lead}(1)
  &=1\cdot1^d\cdot R_{\rm Boc}^{\rm der}(S_E)
   =R_{\rm Boc}^{\rm der}(S_E)\ne0.
 \end{aligned}
\end{equation}
Por tanto el exponente y el coeficiente principal se deducen del germen
holomorfo no nulo; no son entradas de $T_E$, $u_E$ o $S_E$.
Defínanse ahora, después de esta deducción,
\begin{equation}\label{eq:pdf2-bsd-dan-zan}
 d_{\rm an}:=\operatorname{ord}_{s=1}L(E,s),
 \qquad
 \mathfrak z_{\rm an}:=
 \left[(s-1)^{-d_{\rm an}}L(E,s)\right]_{s=1}.
\end{equation}
La ecuación \eqref{eq:pdf2-bsd-expansion-local-deducida} prueba
\begin{equation}\label{eq:pdf2-bsd-dan-d}
 d_{\rm an}=d,
 \qquad
 \mathfrak z_{\rm an}
 =R_{\rm Boc}^{\rm der}(S_E)
 \quad\text{en la línea basada}.
\end{equation}

En el lado aritmético, el teorema~\ref{thm:pdf2-bsd-psi-pt} y la sucesión
\eqref{eq:pdf2-bsd-sha-tamagawa} construyen
\begin{equation}\label{eq:pdf2-bsd-dar-zar}
 \begin{aligned}
 d_{\rm ar}&:=\dim_K M_E^{\rm MW}=d,\\
 \mathfrak z_{\rm ar}&:=
 \frac{|\operatorname{Sha}(E)|\operatorname{Reg}(E)\prod_vc_v}
 {|E(\mathbb Q)_{\rm tors}|^2}
 \quad\text{en la línea aritmética orientada}.
 \end{aligned}
\end{equation}

Para tipar las cuatro flechas, si $C^\bullet$ es un complejo perfecto con
base ordenada $\mathcal B_i=(c_{i,1},\ldots,c_{i,r_i})$, se usa
\begin{equation}\label{eq:pdf2-bsd-linea-determinante-base}
 \det C^\bullet:=
 \bigotimes_i\left(\bigwedge^{r_i}C^i\right)^{(-1)^i},
 \qquad
 \mathbf e_C:=\bigotimes_i
 (c_{i,1}\wedge\cdots\wedge c_{i,r_i})^{(-1)^i}.
\end{equation}
Para un grupo finito $A$ con presentación basada
$\mathbb Z^t\xrightarrow{D_A}\mathbb Z^t\to A$, su línea
$\det_{\mathbb Z}A$ lleva el generador $\mathbf e_A$ y la evaluación
\begin{equation}\label{eq:pdf2-bsd-evaluacion-grupo-finito}
 \operatorname{ev}_A(\mathbf e_A):=|\det D_A|=|A|.
\end{equation}
Cuando esa línea se combina con una línea sobre $K=\mathbb Q$, se usa siempre
la extensión de escala
\begin{equation}\label{eq:pdf2-bsd-linea-finita-k}
 \det_KA:=\det_{\mathbb Z}A\otimes_{\mathbb Z}K;
\end{equation}
no se tensorizan directamente módulos sobre anillos distintos.
En particular, se eligen las bases de Smith ya construidas en
\eqref{eq:pdf2-bsd-smith-finito}, $m_1,\ldots,m_t$ de $M_E$ y
$n_1,\ldots,n_t$ de $N_E$, con $D_Em_i=s_in_i$, y se fija
\begin{equation}\label{eq:pdf2-bsd-generador-smith-fe}
 \mathbf e_{F_E}:=
 (m_1\wedge\cdots\wedge m_t)\otimes
 (n_1\wedge\cdots\wedge n_t)^{-1},
 \qquad
 \operatorname{ev}_{F_E}(\mathbf e_{F_E})=\prod_i|s_i|.
\end{equation}
Las líneas sucesivas son
\begin{equation}\label{eq:pdf2-bsd-cinco-lineas-tipadas}
\begin{aligned}
 \Delta_K&:=\det\mathcal C_E^K,\\
 \Delta_{\rm FERS}&:=\det\mathcal C_E^{\rm Iw},\\
 \Delta_{\rm PT}&:=
   \det_K M_E^{\rm MW}\otimes
   \det_K(M_E^{\rm MW})^\vee\otimes\det_KF_E,\\
 \Delta_{\rm STS}&:=
   \det_K M_E^{\rm MW}\otimes
   \det_K(M_E^{\rm MW})^\vee\otimes
   \det_K\operatorname{Sha}(E)
   \otimes\bigotimes_{v<\infty}
      \det_K\Phi_v(\mathbb F_v),\\
 \Delta_{\rm BSD}&:=\mathbb R\,\mathbf e_{\rm BSD}.
\end{aligned}
\end{equation}
La especialización principal no es una identificación tácita. Denótese
$I:=(T)\subset K[[T]]$ y
$\mathfrak c_{K/{\rm FERS}}^{\rm reg}:=
\det_{\rm alt}\Phi_{K/{\rm FERS}}:\Delta_K[[T]]\to
\Delta_{\rm FERS}[[T]]$. Sean
\begin{equation}\label{eq:pdf2-bsd-lineas-leading}
 \begin{aligned}
  \Delta_{\rm FERS}^{(d)}&:=I^d\Delta_{\rm FERS}[[T]],&
  \Delta_{\rm FERS}^{\rm lead}
    &:=I^{-d}\Delta_{\rm FERS}^{(d)}
       \otimes_{K[[T]],\,T\mapsto0}K,\\
  \Delta_K^{(d)}&:=(\mathfrak c_{K/{\rm FERS}}^{\rm reg})^{-1}
       (\Delta_{\rm FERS}^{(d)}),&
  \Delta_K^{\rm lead}
    &:=I^{-d}\Delta_K^{(d)}
       \otimes_{K[[T]],\,T\mapsto0}K.
 \end{aligned}
\end{equation}
Aquí $I^{-d}\Delta^{(d)}$ significa el módulo formado por los símbolos
$T^{-d}z$ con $z\in\Delta^{(d)}$. La especialización está definida sólo
sobre la parte de orden al menos $d$:
\begin{equation}\label{eq:pdf2-bsd-sp-leading-tipado}
 \operatorname{sp}^{(d)}_{\rm FERS}:
 \Delta_{\rm FERS}^{(d)}\longrightarrow
 \Delta_{\rm FERS}^{\rm lead},
 \qquad
 z(T)\longmapsto
 \overline{T^{-d}z(T)}.
\end{equation}
La equivalencia Kato--FERS induce entonces el mapa tipado
\begin{equation}\label{eq:pdf2-bsd-c-k-fers-leading}
 \begin{aligned}
 \mathfrak c_{K/{\rm FERS}}^{\rm lead}:{}
 &\Delta_K^{\rm lead}\xrightarrow{\ \sim\ }
   \Delta_{\rm FERS}^{\rm lead},\\
 &\overline{T^{-d}z_K(T)}\longmapsto
   \overline{T^{-d}
   \mathfrak c_{K/{\rm FERS}}^{\rm reg}(z_K(T))}.
 \end{aligned}
\end{equation}
La fórmula no aplica $\operatorname{sp}^{(d)}$ a una línea desnuda. En
particular, si
\begin{equation}\label{eq:pdf2-bsd-seccion-zeta-leading}
 \begin{aligned}
  \boldsymbol\zeta_{\rm FERS}(T)
    &:=\det S_E(T)\,\mathbf e_{\rm FERS}
       \in\Delta_{\rm FERS}^{(d)},\\
  \boldsymbol\zeta_K(T)
    &:=(\mathfrak c_{K/{\rm FERS}}^{\rm reg})^{-1}
       \boldsymbol\zeta_{\rm FERS}(T)
       \in\Delta_K^{(d)},
 \end{aligned}
\end{equation}
entonces
\begin{equation}\label{eq:pdf2-bsd-especializacion-zeta}
 \operatorname{sp}^{(d)}_{\rm FERS}
       (\boldsymbol\zeta_{\rm FERS})
 =R_{\rm Boc}^{\rm der}(S_E)\,
       \mathbf e_{\rm FERS}^{\rm lead},
 \quad
 \mathfrak c_{K/{\rm FERS}}^{\rm lead}
       (\overline{T^{-d}\boldsymbol\zeta_K})
 =R_{\rm Boc}^{\rm der}(S_E)\,
       \mathbf e_{\rm FERS}^{\rm lead}.
\end{equation}
La definición usa sólo $d=\operatorname{ord}_T\det S_E(T)$, ya calculado por
Smith, y no el orden analítico.

Las cuatro flechas propietarias, ahora con dominio y codominio explícitos,
se comparan en la fibra real desde el primer lugar donde interviene la
altura de Néron.  Para evitar multiplicar una línea sobre $K$ por un
escalar real no perteneciente en general a $K$, se fijan
\begin{equation}\label{eq:pdf2-bsd-lineas-reales-comparador}
 \begin{aligned}
  \Delta_{K,\mathbb R}^{\rm lead}
    &:=\Delta_K^{\rm lead}\otimes_{K,\iota_\infty}\mathbb R,&
  \Delta_{{\rm FERS},\mathbb R}^{\rm lead}
    &:=\Delta_{\rm FERS}^{\rm lead}\otimes_{K,\iota_\infty}\mathbb R,\\
  \Delta_{{\rm PT},\mathbb R}
    &:=\Delta_{\rm PT}\otimes_{K,\iota_\infty}\mathbb R,&
  \Delta_{{\rm STS},\mathbb R}
    &:=\Delta_{\rm STS}\otimes_{K,\iota_\infty}\mathbb R.
 \end{aligned}
\end{equation}
Las cuatro flechas son entonces
\begin{equation}\label{eq:pdf2-bsd-cuatro-comparadores}
\begin{aligned}
 \mathfrak c_{K/{\rm FERS},\mathbb R}^{\rm lead}:
   &\Delta_{K,\mathbb R}^{\rm lead}
      \longrightarrow\Delta_{{\rm FERS},\mathbb R}^{\rm lead},
 &\mathfrak c_{K/{\rm FERS},\mathbb R}^{\rm lead}
   &:=\mathfrak c_{K/{\rm FERS}}^{\rm lead}\otimes_K1_{\mathbb R},\\
 \mathfrak c_{{\rm PT},\mathbb R}:
   &\Delta_{{\rm FERS},\mathbb R}^{\rm lead}
      \longrightarrow\Delta_{{\rm PT},\mathbb R},
 &\mathfrak c_{{\rm PT},\mathbb R}(z)&:=
   \mathfrak n_{E,\mathbb R}(z)\otimes\mathbf e_{F_E},\\
 \mathfrak c_{{\rm STS},\mathbb R}:
   &\Delta_{{\rm PT},\mathbb R}
      \longrightarrow\Delta_{{\rm STS},\mathbb R},
 &\mathfrak c_{{\rm STS},\mathbb R}&:=\det\bigl(
   0\to\operatorname{Sha}(E)\to F_E
     \to\textstyle\bigoplus_{v<\infty}\Phi_v(\mathbb F_v)
     \to0\bigr)\otimes_K1_{\mathbb R},\\
 \mathfrak c_{\rm tors,\infty}:
   &\Delta_{{\rm STS},\mathbb R}
      \longrightarrow\Delta_{\rm BSD}.&&
\end{aligned}
\end{equation}

\begin{proposition}[Cálculo generador a generador de las cuatro flechas]
\label{prop:pdf2-bsd-cuatro-flechas-generadores}
Las acciones sobre las bases de
\eqref{eq:pdf2-bsd-cinco-lineas-tipadas} son las siguientes.

\emph{Primera flecha.} Si
$\mathbf e_K^i=\bigwedge_{b\in\mathcal B_i^K}b$ y
$\mathbf e_{\rm FERS}^i=\bigwedge_{b\in\mathcal B_i^K}\upsilon_i(b)$,
la tabla \eqref{eq:pdf2-bsd-tabla-kato-fers} da
\begin{equation}\label{eq:pdf2-bsd-c1-generadores}
 \mathfrak c_{K/{\rm FERS}}^{\rm reg}(\mathbf e_K)
 =\prod_i\det(I+TB_i(T))^{(-1)^i}\mathbf e_{\rm FERS}
 =\rho_E(T)\mathbf e_{\rm FERS},
 \qquad \rho_E(0)=1.
\end{equation}
Si $\mathbf e_K^{\rm lead}$ y $\mathbf e_{\rm FERS}^{\rm lead}$ son las
bases inducidas por \eqref{eq:pdf2-bsd-lineas-leading}, entonces
\begin{equation}\label{eq:pdf2-bsd-c1-leading-generador}
 \mathfrak c_{K/{\rm FERS},\mathbb R}^{\rm lead}
     (\mathbf e_K^{\rm lead}\otimes1)
 =\rho_E(0)(\mathbf e_{\rm FERS}^{\rm lead}\otimes1)
 =\mathbf e_{\rm FERS}^{\rm lead}\otimes1.
\end{equation}

\emph{Segunda flecha.} Elíjanse bases integrales
$a_{q,1},\ldots,a_{q,r_q}$ de $A_{q,\mathbb Z}$, vistas en $A_q$, y
$b_{q,j}=\bar\beta_q(a_{q,j})$ de $B_q$. La matriz de la altura de la página
queda definida por
\begin{equation}\label{eq:pdf2-bsd-hq-generadores}
 j_q(b_{q,j})=
 \sum_{k=1}^{r_q}(H_q)_{kj}\,a_{q,k}^\vee\otimes T^q.
\end{equation}
Por tanto, sobre la cuña completa,
\begin{equation}\label{eq:pdf2-bsd-c2-cuna-paginas}
 \bigwedge_{j=1}^{r_q}j_q\bar\beta_q(a_{q,j})
 =T^{qr_q}\det(H_q)
  \bigwedge_{k=1}^{r_q}a_{q,k}^\vee.
\end{equation}
El producto sobre $q$, junto con las filas regulares, tiene exponente
$\sum_q qr_q=\sum_q\dim_K V_q=d$. Para
$1\leq j\leq r_q$ y $1\leq k\leq q$ póngase
$e_{q,j,k}:=a_{q,j}e_{q,k}$ en
$\operatorname{Jet}_{\rm Boc,\mathbb Z}(E)$ y defínase
\begin{equation}\label{eq:pdf2-bsd-base-mw-desde-jets}
 P_{q,j,k}:=\Psi_{{\rm PT},\mathbb Z}^{-1}
       \mathcal U_{\rm jet}(e_{q,j,k})
 \in M_{E,\mathbb Z}^{\rm MW}.
\end{equation}
Por \eqref{eq:pdf2-bsd-desdoblamiento-jets},
\eqref{eq:pdf2-bsd-dimension-jet} y
\eqref{eq:pdf2-bsd-psi-pt-integral}, estos $d$ elementos son una base
integral, no una subred de índice no controlado. Se enumeran en orden
lexicográfico como $P_1,\ldots,P_d$ y, sobre $K$, forman la base de
$M_E^{\rm MW}$ usada a continuación.

La altura se construye independientemente de FERS.\@ Para
$P\in E(\overline{\mathbb Q})$ se fija
\begin{equation}\label{eq:pdf2-bsd-altura-canonica-construida}
 \widehat h_E(P):=
 \frac12\lim_{n\to\infty}4^{-n}h_x([2^n]P),
 \qquad
 \langle P,Q\rangle_{\rm NT}:=
 \frac12\bigl(\widehat h_E(P+Q)-\widehat h_E(P)-\widehat h_E(Q)\bigr).
\end{equation}
Equivalentemente, el biextensor de Poincaré rigidificado sobre el modelo de
Néron produce, para todo lugar $v$, incluido $v=\infty$, el símbolo local
real $\lambda_v(P,Q)$ normalizado por la sección cero, y
\begin{equation}\label{eq:pdf2-bsd-altura-neron-local-global}
 \langle P,Q\rangle_{\rm NT}
 =\sum_v\lambda_v(P,Q)\in\mathbb R.
\end{equation}
No se sustituye esta altura real, simétrica y positiva por una suma de
invariantes de cup-products con valores en $\mathbb Q/\mathbb Z$.

Póngase
$H_{\rm NT}^{\flat}(P)(Q):=\langle P,Q\rangle_{\rm NT}$ y sea
$\mathcal B_E:=\Psi_{{\rm PT},\mathbb Z}^{-1}\mathcal U_{\rm jet}$,
extendido a $\mathbb R$. No se sustituyen las $q-1$ direcciones adicionales
por identidades: se calcula el bloque Gram completo, incluidos todos los
cruces entre profundidades, filas y jets,
\begin{equation}\label{eq:pdf2-bsd-gram-jet-completo}
 \begin{aligned}
 G_E^{\rm jet}\bigl((q,j,k),(q',j',k')\bigr)
 &:={}
 \left\langle P_{q,j,k},P_{q',j',k'}\right\rangle_{\rm NT}\\
 &=\sum_v\lambda_v
   \left(P_{q,j,k},P_{q',j',k'}\right),
 \end{aligned}
\end{equation}
para todo $1\le k\le q$, $1\le k'\le q'$. Esta matriz es material: los
$P_{q,j,k}$ se construyeron en
\eqref{eq:pdf2-bsd-base-mw-desde-jets} y cada entrada se calcula por
\eqref{eq:pdf2-bsd-altura-neron-local-global}. Por definición del pullback
de una forma bilineal,
\begin{equation}\label{eq:pdf2-bsd-altura-nt-generadores}
 \boxed{
 \mathcal B_E^\vee H_{\rm NT}^{\flat}\mathcal B_E=G_E^{\rm jet}.}
\end{equation}
No se afirma que $G_E^{\rm jet}$ sea diagonal por bloques. En la base
integral $P_1,\ldots,P_d$, el regulador y su acción en la línea determinante
son exactamente
\begin{equation}\label{eq:pdf2-bsd-regulador-cuna-altura}
 \begin{aligned}
 \operatorname{Reg}(E)&:=\det G_E^{\rm jet}>0,\\
 \bigwedge_{j=1}^dH_{\rm NT}^{\flat}(P_j)
 &=\operatorname{Reg}(E)
   (P_1^\vee\wedge\cdots\wedge P_d^\vee).
 \end{aligned}
\end{equation}
La página Bockstein y la altura son dos emparejamientos construidos, con
bases ya fijadas. El comparador de Néron entre sus líneas se define antes de
usar el término principal por
\begin{equation}\label{eq:pdf2-bsd-comparador-neron-jet}
 \begin{aligned}
 \mathfrak n_{E,\mathbb R}:
 \Delta_{{\rm FERS},\mathbb R}^{\rm lead}&\longrightarrow
 \bigl(\det_K(M_E^{\rm MW})\otimes_K
 \det_K(M_E^{\rm MW})^\vee\bigr)\otimes_{K,\iota_\infty}\mathbb R,\\
 \mathfrak n_{E,\mathbb R}
 (\mathbf e_{\rm FERS}^{\rm lead}\otimes1)
 &:=\frac{\operatorname{Reg}(E)}{R_{\rm Boc}^{\rm der}(S_E)}
 (P_1\wedge\cdots\wedge P_d)
 \otimes(P_1^\vee\wedge\cdots\wedge P_d^\vee)\otimes1.
 \end{aligned}
\end{equation}
El denominador es no nulo por \eqref{eq:pdf2-bsd-a-e-cero}; numerador y
denominador se obtienen, respectivamente, del Gram completo
\eqref{eq:pdf2-bsd-gram-jet-completo} y de las páginas
\eqref{eq:pdf2-bsd-rboc-lt}. Así la normalización es un cociente real de dos
determinantes previamente calculados, no una ortogonalidad inventada ni una
consulta de la fórmula BSD. La extensión de escala se efectúa antes de
aplicar el regulador y se conserva en las flechas PT y STS.

Las páginas superiores y la fila regular conservan el generador integral de
$F_E$ mediante un mapa tipado. Se usa la extensión racional de la sección
$s_q$ ya construida en \eqref{eq:pdf2-bsd-desdoblamiento-jets}. La sección dual
\begin{equation}\label{eq:pdf2-bsd-glue-dual-tipado}
 \operatorname{gl}_q^\vee:A_q^\vee\longrightarrow
   (\operatorname{Surv}_{\rm det}^{\rm MW})^\vee
\end{equation}
se define en la base completa por
$(\operatorname{gl}_q^\vee\varphi)
(\operatorname{gl}_r s_r(a))=\delta_{qr}\varphi(a)$ y por cero en las filas
FS.\@ Sea $\mathcal I_{\rm fin}$ la lista ordenada de filas regulares y pares
$(q,j)$ con $q\geq2$. Para $\eta$ con $q(\eta)\geq2$ se ponen
\begin{equation}\label{eq:pdf2-bsd-lifts-filas-smith}
 \begin{aligned}
 \widetilde a_\eta&:=\operatorname{pr}_{M_E}\operatorname{pr}_{\mathcal L_E}
   \Theta_E\operatorname{gl}_{q(\eta)}s_{q(\eta)}(a_\eta),\\
 \widetilde b_\eta^\vee&:=
   \operatorname{pr}_{N_E^\vee}\operatorname{pr}_{\mathcal L_E^\vee}
   (\Theta_E^\vee)^{-1}\operatorname{gl}_{q(\eta)}^\vee
   \bigl((1\otimes\partial_T^{(q(\eta))})
      j_{q(\eta)}\bar\beta_{q(\eta)}(a_\eta)\bigr).
 \end{aligned}
\end{equation}
En una fila regular se definen separadamente
\begin{equation}\label{eq:pdf2-bsd-lifts-fila-regular}
 \widetilde a_\eta:=
 \operatorname{pr}_{M_E}\operatorname{pr}_{\mathcal L_E}
 \Theta_Es_{\rm reg}(a_\eta),
 \qquad
 \widetilde b_\eta^\vee:=
 \operatorname{pr}_{N_E^\vee}\operatorname{pr}_{\mathcal L_E^\vee}
 (\Theta_E^\vee)^{-1}s_{\rm reg}^\vee
 (\bar S_E(0)a_\eta).
\end{equation}
Todos los dominios son ahora $A_q$ o $A_q^\vee$ según corresponde; no se
aplica $\operatorname{gl}_q$ a un covector.

La biyección requerida no se infiere de matrices unimodulares. Las
ecuaciones de unidad--counidad muestran sobre cada generador que las listas
$(\widetilde a_\eta)_\eta$ y sus terminales duales son bases de $M_E$ y
$N_E^\vee$. Escríbase la lista heredada como
$\mathcal I_{\rm fin}=\{\eta_1<\cdots<\eta_t\}$ y defínase directamente
\begin{equation}\label{eq:pdf2-bsd-nu-construida}
 \boxed{\nu(\eta_j):=j},
 \qquad
 \nu:\mathcal I_{\rm fin}\xrightarrow{\sim}\{1,\ldots,t\}.
\end{equation}
Pónganse $\mu_j:=\widetilde a_{\eta_j}$ y sea
$\xi_1,\ldots,\xi_t$ la base terminal de $N_E$. Antes de Smith, la matriz
material es
\begin{equation}\label{eq:pdf2-bsd-c2-filas-finitas}
 D^{\rm term}_{kj}:=
 \langle\xi_k^\vee,D_E\mu_j\rangle,
 \qquad
 \widetilde b_{\eta_j}^\vee
 =\sum_{k=1}^tD^{\rm term}_{kj}\xi_k^\vee.
\end{equation}
Esta igualdad es la relación terminal de
\eqref{eq:pdf2-bsd-descomposicion-global}, evaluada en la counidad
$\xi_k^\vee$; no divide por un divisor elemental. Si
$UD^{\rm term}V=\operatorname{diag}(s_1,\ldots,s_t)$, entonces
\begin{equation}\label{eq:pdf2-bsd-cuna-smith-con-uv}
 \bigwedge_{j=1}^t\widetilde b_{\eta_j}^\vee
 =\det(D^{\rm term})
   (\xi_1^\vee\wedge\cdots\wedge\xi_t^\vee),
 \qquad
 |\det(D^{\rm term})|=\prod_{j=1}^t|s_j|.
\end{equation}
Los factores $\det U,\det V\in\{\pm1\}$ quedan registrados por
$\operatorname{or}_{\rm det}$ al pasar de $(\mu_j),(\xi_j)$ a las bases de
Smith $(m_j),(n_j)$. Por tanto el cociente de cuñas es exactamente
$\mathbf e_{F_E}$ de \eqref{eq:pdf2-bsd-generador-smith-fe}; no se ha
supuesto que $U$ o $V$ sean matrices de permutación y no queda una página
sin publicar. En
consecuencia,
\begin{equation}\label{eq:pdf2-bsd-c2-generadores}
 \mathfrak c_{{\rm PT},\mathbb R}
 \bigl(R_{\rm Boc}^{\rm der}(S_E)
       (\mathbf e_{\rm FERS}^{\rm lead}\otimes1)\bigr)
 =\operatorname{Reg}(E)\,
  (P_1\wedge\cdots\wedge P_d)\otimes
  (P_1^\vee\wedge\cdots\wedge P_d^\vee)
  \otimes\mathbf e_{F_E}\otimes1.
\end{equation}

\emph{Tercera flecha.} Para cada $\ell$, elíjanse generadores ordenados
$\sigma_{\ell,a}$ de $\operatorname{Sha}(E)[\ell^\infty]$ y
$\phi_{\ell,v,b}$ de $\Phi_v(\mathbb F_v)[\ell^\infty]$. Los levantamientos explícitos
$\operatorname{Lift}_{\Phi,\ell}(\phi_{\ell,v,b})$ y las imágenes
$\iota_{{\rm Sha},\ell}(\sigma_{\ell,a})$ forman, por exactitud, una base de
presentación de $F_{E,\ell}$. En esas bases la matriz de relaciones se
calcula generador a generador como
\begin{equation}\label{eq:pdf2-bsd-c3-matriz-presentacion}
 D_{F,\ell}=
 \begin{pmatrix}
  D_{{\rm Sha},\ell}&X_\ell\\
  0&\displaystyle\bigoplus_{v<\infty}D_{\Phi_v,\ell}
 \end{pmatrix}.
\end{equation}
La columna $X_\ell$ registra el cambio de cada levantamiento al modificar su
representante; el bloque inferior izquierdo es cero porque
$\partial_{{\rm loc},\ell}\iota_{{\rm Sha},\ell}=0$. Por tanto la cuña de
presentación no depende de $X_\ell$ y
$\det D_{F,\ell}=\det D_{{\rm Sha},\ell}
\prod_{v<\infty}\det D_{\Phi_v,\ell}$. Sobre ella,
\begin{equation}\label{eq:pdf2-bsd-c3-generadores}
 \mathbf e_{F_{E,\ell}}
 \longmapsto
 \mathbf e_{\operatorname{Sha}(E)[\ell^\infty]}
 \otimes\bigotimes_{v<\infty}
    \mathbf e_{\Phi_v(\mathbb F_v)[\ell^\infty]},
\end{equation}
y las evaluaciones de \eqref{eq:pdf2-bsd-evaluacion-grupo-finito} dan
$|F_{E,\ell}|=|\operatorname{Sha}(E)[\ell^\infty]|
\prod_{v<\infty}|\Phi_v(\mathbb F_v)[\ell^\infty]|$.
Multiplicar sobre $\ell$ produce
\begin{equation}\label{eq:pdf2-bsd-c3-integral-generadores}
 \mathfrak c_{{\rm STS},\mathbb R}(\mathbf e_{F_E}\otimes1)
 =(\mathbf e_{\operatorname{Sha}(E)}
  \otimes\bigotimes_{v<\infty}\mathbf e_{\Phi_v(\mathbb F_v)})\otimes1,
 \qquad
 |F_E|=|\operatorname{Sha}(E)|\prod_vc_v.
\end{equation}

\emph{Cuarta flecha.} Fijada la base de Mordell--Weil anterior y una
diferencial de Néron $\omega_E$ con período positivo
$\Omega_E$, la base orientada de $\Delta_{\rm STS}$ se envía por
\begin{equation}\label{eq:pdf2-bsd-c4-generadores}
 \begin{aligned}
 &x\,(P_1\wedge\cdots\wedge P_d)\otimes
 (P_1^\vee\wedge\cdots\wedge P_d^\vee)\otimes
 \mathbf e_{\operatorname{Sha}(E)}
 \otimes\bigotimes_{v<\infty}\mathbf e_{\Phi_v(\mathbb F_v)}\\
 &\hspace{28mm}\longmapsto
 \frac{x\,\Omega_E\,
   \operatorname{ev}_{\operatorname{Sha}(E)}
      (\mathbf e_{\operatorname{Sha}(E)})
   \prod_{v<\infty}\operatorname{ev}_{\Phi_v(\mathbb F_v)}
      (\mathbf e_{\Phi_v(\mathbb F_v)})}
 {|E(\mathbb Q)_{\rm tors}|^2}\,
 \mathbf e_{\rm BSD}.
\end{aligned}
\end{equation}
Por \eqref{eq:pdf2-bsd-evaluacion-grupo-finito}, el coeficiente de la última
línea es
$x\Omega_E|\operatorname{Sha}(E)|\prod_vc_v/
|E(\mathbb Q)_{\rm tors}|^2$.
Las dos copias del grupo de torsión proceden de los grados duales $0$ y $2$;
el factor $\Omega_E$ es la evaluación de $\omega_E$ en la componente real
orientada. Esta orientación coloca $\Omega_E\mathfrak z_{\rm ar}$, y no su
inverso, como coeficiente aritmético de la composición; la comparación con
$\mathfrak z_{\rm an}$ se deducirá sólo después mediante el cuadrado de
trivializaciones.
\end{proposition}

\begin{proof}
La primera identidad es el determinante alternado de las matrices triangulares
$I+TB_i(T)$. La segunda aplica la especialización de las páginas y después
el comparador de Néron \eqref{eq:pdf2-bsd-comparador-neron-jet}. Su
numerador es el determinante del Gram jet completo
\eqref{eq:pdf2-bsd-gram-jet-completo}, con todos los cruces, y su denominador
es el determinante Bockstein ya calculado. Por
\eqref{eq:pdf2-bsd-regulador-cuna-altura} su acción sobre cuñas es
\eqref{eq:pdf2-bsd-c2-generadores}; el exponente $d$ se obtiene antes de
tomar dimensiones de Mordell--Weil. La tercera es el isomorfismo canónico de líneas asociado a la
sucesión exacta, escrito en las secciones
$\iota_{{\rm Sha},\ell}$ y $\operatorname{Lift}_{\Phi,\ell}$ ya
construidas. La cuarta es la evaluación de los dos complejos finitos de
torsión y de la línea de Néron. Esto calcula las cuatro flechas sin usar la
igualdad BSD.
\end{proof}

Su composición
\begin{equation}\label{eq:pdf2-bsd-comparador-total}
 \mathfrak C_E:=
 \mathfrak c_{\rm tors,\infty}\circ
 \mathfrak c_{{\rm STS},\mathbb R}\circ
 \mathfrak c_{{\rm PT},\mathbb R}\circ
 \mathfrak c_{K/{\rm FERS},\mathbb R}^{\rm lead}:
 \Delta_{K,\mathbb R}^{\rm lead}
 \longrightarrow\Delta_{\rm BSD}
\end{equation}
queda bien tipada. Antes de insertar el elemento analítico o cardinal alguno
se fijan las trivializaciones independientes
\begin{equation}\label{eq:pdf2-bsd-trivializaciones-independientes}
 \operatorname{tr}_{\rm an}(x\mathbf e_K^{\rm lead}):=x,
 \qquad
 \operatorname{tr}_{\rm Ner}(y\mathbf e_{\rm BSD}):=y.
\end{equation}
La primera usa sólo la base Kato especializada por
\eqref{eq:pdf2-bsd-lineas-leading} y la normalización Mellin--Fredholm de
\eqref{eq:pdf2-bsd-fredholm-euler}; en particular,
$\operatorname{tr}_{\rm an}
(\overline{T^{-d}\boldsymbol\zeta_K})
=R_{\rm Boc}^{\rm der}(S_E)=\mathfrak z_{\rm an}$ por
\eqref{eq:pdf2-bsd-especializacion-zeta}. La segunda usa sólo la diferencial
de Néron orientada y satisface
$\operatorname{tr}_{\rm Ner}(\mathbf e_{\rm BSD})=1$. La identidad del
comparador basado de $\mathfrak G_{\rm det}(E)$, desplegada por las cuatro
flechas, es el
cuadrado
\begin{equation}\label{eq:pdf2-bsd-cuadrado-trivializaciones}
\begin{array}{ccc}
 \Delta_{K,\mathbb R}^{\rm lead}
   &\xrightarrow{\ \mathfrak C_E\ }&\Delta_{\rm BSD}\\[2mm]
 \big\downarrow\scriptstyle\operatorname{tr}_{\rm an}
   &&\big\downarrow\scriptstyle\operatorname{tr}_{\rm Ner}\\[2mm]
 \mathbb R&=&\mathbb R.
\end{array}
\end{equation}
Este cuadrado se fija con las bases estructurales antes de evaluar
$L(E,s)$, el regulador o los grupos finitos.

El elemento analítico se coloca, después de construir las flechas, en su
dominio correcto:
\begin{equation}\label{eq:pdf2-bsd-elemento-analitico-tipado}
 \mathbf z_E^{\rm an}:=\mathfrak z_{\rm an}\mathbf e_K^{\rm lead}
 =\overline{T^{-d}\boldsymbol\zeta_K}
 \in\Delta_{K,\mathbb R}^{\rm lead}.
\end{equation}
Por \eqref{eq:pdf2-bsd-dan-d},
$\mathfrak z_{\rm an}=R_{\rm Boc}^{\rm der}(S_E)$; la segunda igualdad de
\eqref{eq:pdf2-bsd-especializacion-zeta} prueba que la primera flecha lleva
$\mathbf z_E^{\rm an}$ exactamente al miembro izquierdo de
\eqref{eq:pdf2-bsd-c2-generadores}. Las otras tres flechas, calculadas sobre
generadores, dan
\begin{equation}\label{eq:pdf2-bsd-elementos-linea}
 d_{\rm an}=d=d_{\rm ar},
 \qquad
 \boxed{
 \mathfrak C_E(\mathbf z_E^{\rm an})
 =\Omega_E\mathfrak z_{\rm ar}\,\mathbf e_{\rm BSD}.}
\end{equation}
Aplicar las flechas verticales de
\eqref{eq:pdf2-bsd-cuadrado-trivializaciones} produce la igualdad numérica
\begin{equation}\label{eq:pdf2-bsd-igualdad-numerica-deducida}
 \boxed{\mathfrak z_{\rm an}=\Omega_E\mathfrak z_{\rm ar}.}
\end{equation}

\begin{theorem}[Publicación de Birch--Swinnerton--Dyer]
\label{thm:pdf2-bsd-publicacion-final}
Para la imagen determinantal autónoma completa
$\mathfrak G_{\rm det}(E)$ se tiene
\begin{equation}\label{eq:pdf2-bsd-rango-final}
 \operatorname{ord}_{s=1}L(E,s)=\operatorname{rank}E(\mathbb Q)=:r
\end{equation}
y
\begin{equation}\label{eq:pdf2-bsd-formula-final}
 \frac{L^{(r)}(E,1)}{r!\,\Omega_E}
 =\frac{|\operatorname{Sha}(E)|\operatorname{Reg}(E)\prod_vc_v}
 {|E(\mathbb Q)_{\rm tors}|^2}.
\end{equation}
\end{theorem}

\begin{proof}
La igualdad de grados es
\eqref{eq:pdf2-bsd-dan-d} seguida de
\eqref{eq:pdf2-bsd-rango-orden}. La tabla Kato--FERS transporta el elemento
analítico a todas las páginas Bockstein; $\Psi_{\rm PT}$ las pega a la red
de Mordell--Weil y publica su determinante de altura; la sucesión exacta
Sha--Tamagawa descompone el factor de Smith; y el último comparador inserta
la red torsional y el período ya orientado. Esto es exactamente
\eqref{eq:pdf2-bsd-elementos-linea}. El cuadrado independiente
\eqref{eq:pdf2-bsd-cuadrado-trivializaciones} da
\eqref{eq:pdf2-bsd-igualdad-numerica-deducida}; sustituir las definiciones de
$\mathfrak z_{\rm an}$ y $\mathfrak z_{\rm ar}$ produce
\eqref{eq:pdf2-bsd-formula-final}. Ninguna de las cuatro flechas se
normaliza con $r$, $\operatorname{Reg}(E)$,
$|\operatorname{Sha}(E)|$ o el término principal: esos valores aparecen
sólo al final de sus respectivos determinantes.
\end{proof}

\section{Controles no gobernantes, falsadores y estatuto}

\begin{falsifier}[Falsadores separados por componente]
La construcción falla si una historia de
$\Sigma_{\rm det}^{\rm mult}$ no pertenece a la partición
\eqref{eq:pdf2-bsd-tres-proyectores}; si se escribe
$\partial H_{\rm FS}+H_{\rm FS}\partial=q_{\rm root}$ y con ello se borra
$\mathcal L_E$; si una fila no raíz de
\eqref{eq:pdf2-bsd-tabla-kato-fers} tiene peso cero; si fallan las ecuaciones
unidad--counidad \eqref{eq:pdf2-bsd-unidad-counidad-paginas}; si
$D_E\otimes\mathbb Q$ no es invertible; si
$\pi_\Phi\operatorname{Lift}_\Phi\ne I$; o si alguna flecha de
\eqref{eq:pdf2-bsd-cuatro-comparadores} usa el valor final para fijar su
base. Cada falsador rompe una arista concreta y no reabre las anteriores.
\end{falsifier}

La homotopía $H_{\rm FS}$ es un control posterior del filler explícito
$h_*=-vf$ y no una premisa gobernante. Las formas de Smith, los
determinantes y los cardinales se calculan después de construir los mapas;
no sustituyen las tablas, el pegado de páginas ni la exactitud de cociclos.

\paragraph{Estatuto y procedencia.}
Alexander--Whitney, el producto fibrado, la unidad común, la dualidad
reducida, la rigidez de la raíz, el filler, todas las páginas Bockstein y
los cuatro lectores pertenecen a
\texttt{ARQUITECTURA\_AUTORAL\_PREEXISTENTE}/\texttt{EXACTO\_INTERNO} de
$\mathfrak G_{\rm det}(E)$. La partición exhaustiva
\eqref{eq:pdf2-bsd-descomposicion-global}, la tabla
\eqref{eq:pdf2-bsd-tabla-kato-fers}, las fórmulas
\eqref{eq:pdf2-bsd-psi-y-lambda-q} y los mapas de cociclos
\eqref{eq:pdf2-bsd-iota-sha-cociclos}--
\eqref{eq:pdf2-bsd-lift-componentes} son
\texttt{FORMALIZACION\_NUEVA}: despliegan materialmente los entrelazadores
ya contenidos en las trece coordenadas estructurales. El teorema final tiene
fuerza
\texttt{DEMOSTRADO\_CON\_ESTRUCTURA\_DE\_PARTIDA\_EXPLICITA}; no queda una
obligación residual gobernante ni se usa una homotopía abstracta para
fabricar el resultado.
